\documentclass{ndss}
\usepackage{algorithm,algpseudocode,pifont,subcaption,multirow,xcolor,colortbl}
\newcommand{\cmark}{\ding{51}}
\newcommand{\xmark}{\ding{55}}
\newif\ifanonymous
\anonymoustrue
% shared/macros.tex — Identity-masked macros for both arXiv and NDSS versions
% Toggle anonymization: set \anonymoustrue before % shared/macros.tex — Identity-masked macros for both arXiv and NDSS versions
% Toggle anonymization: set \anonymoustrue before % shared/macros.tex — Identity-masked macros for both arXiv and NDSS versions
% Toggle anonymization: set \anonymoustrue before \input{macros} for NDSS,
%                       set \anonymousfalse before \input{macros} for arXiv

\ifanonymous
  \newcommand{\toolname}{{\sc BoundaryFlow}\xspace}
  \newcommand{\anonrepo}{\url{https://anonymous.4open.science/r/XXXX}}
  \newcommand{\authorblock}{%
    \author{Anonymous Submission \\ Paper ID: TBD}%
  }
\else
  \newcommand{\toolname}{{\sc kcov-dataflow}\xspace}
  \newcommand{\anonrepo}{\url{https://github.com/yskzalloc/kcov-dataflow}}
  \newcommand{\authorblock}{%
    \author{Yunseong Kim \\
      Ericsson Software Technology \\
      \texttt{yunseong.kim@est.tech}}%
    % \author{Anonymous Submission}%
  }
\fi

% Technical shorthands
\newcommand{\sancov}{\texttt{SanitizerCoverage}\xspace}
\newcommand{\tracepc}{\texttt{trace-pc}\xspace}
\newcommand{\traceargs}{\texttt{trace-args}\xspace}
\newcommand{\traceret}{\texttt{trace-ret}\xspace}
\newcommand{\kcov}{\texttt{KCOV}\xspace}
\newcommand{\kasan}{\texttt{KASAN}\xspace}
\newcommand{\syzkaller}{\texttt{syzkaller}\xspace}
\newcommand{\drgn}{\texttt{drgn}\xspace}
\newcommand{\virtme}{\texttt{virtme-ng}\xspace}

% Data-flow tuple notation
\newcommand{\dftuple}[4]{\ensuremath{\langle #1, #2, #3, #4 \rangle}}
\newcommand{\callerid}{\ensuremath{\mathit{PC}_{\mathrm{caller}}}}
\newcommand{\calleeid}{\ensuremath{\mathit{PC}_{\mathrm{callee}}}}

% Listing style
\usepackage{listings}
\lstset{
  basicstyle=\ttfamily\footnotesize,
  breaklines=true,
  frame=single,
  numbers=left,
  numberstyle=\tiny,
  captionpos=b,
}
 for NDSS,
%                       set \anonymousfalse before % shared/macros.tex — Identity-masked macros for both arXiv and NDSS versions
% Toggle anonymization: set \anonymoustrue before \input{macros} for NDSS,
%                       set \anonymousfalse before \input{macros} for arXiv

\ifanonymous
  \newcommand{\toolname}{{\sc BoundaryFlow}\xspace}
  \newcommand{\anonrepo}{\url{https://anonymous.4open.science/r/XXXX}}
  \newcommand{\authorblock}{%
    \author{Anonymous Submission \\ Paper ID: TBD}%
  }
\else
  \newcommand{\toolname}{{\sc kcov-dataflow}\xspace}
  \newcommand{\anonrepo}{\url{https://github.com/yskzalloc/kcov-dataflow}}
  \newcommand{\authorblock}{%
    \author{Yunseong Kim \\
      Ericsson Software Technology \\
      \texttt{yunseong.kim@est.tech}}%
    % \author{Anonymous Submission}%
  }
\fi

% Technical shorthands
\newcommand{\sancov}{\texttt{SanitizerCoverage}\xspace}
\newcommand{\tracepc}{\texttt{trace-pc}\xspace}
\newcommand{\traceargs}{\texttt{trace-args}\xspace}
\newcommand{\traceret}{\texttt{trace-ret}\xspace}
\newcommand{\kcov}{\texttt{KCOV}\xspace}
\newcommand{\kasan}{\texttt{KASAN}\xspace}
\newcommand{\syzkaller}{\texttt{syzkaller}\xspace}
\newcommand{\drgn}{\texttt{drgn}\xspace}
\newcommand{\virtme}{\texttt{virtme-ng}\xspace}

% Data-flow tuple notation
\newcommand{\dftuple}[4]{\ensuremath{\langle #1, #2, #3, #4 \rangle}}
\newcommand{\callerid}{\ensuremath{\mathit{PC}_{\mathrm{caller}}}}
\newcommand{\calleeid}{\ensuremath{\mathit{PC}_{\mathrm{callee}}}}

% Listing style
\usepackage{listings}
\lstset{
  basicstyle=\ttfamily\footnotesize,
  breaklines=true,
  frame=single,
  numbers=left,
  numberstyle=\tiny,
  captionpos=b,
}
 for arXiv

\ifanonymous
  \newcommand{\toolname}{{\sc BoundaryFlow}\xspace}
  \newcommand{\anonrepo}{\url{https://anonymous.4open.science/r/XXXX}}
  \newcommand{\authorblock}{%
    \author{Anonymous Submission \\ Paper ID: TBD}%
  }
\else
  \newcommand{\toolname}{{\sc kcov-dataflow}\xspace}
  \newcommand{\anonrepo}{\url{https://github.com/yskzalloc/kcov-dataflow}}
  \newcommand{\authorblock}{%
    \author{Yunseong Kim \\
      Ericsson Software Technology \\
      \texttt{yunseong.kim@est.tech}}%
    % \author{Anonymous Submission}%
  }
\fi

% Technical shorthands
\newcommand{\sancov}{\texttt{SanitizerCoverage}\xspace}
\newcommand{\tracepc}{\texttt{trace-pc}\xspace}
\newcommand{\traceargs}{\texttt{trace-args}\xspace}
\newcommand{\traceret}{\texttt{trace-ret}\xspace}
\newcommand{\kcov}{\texttt{KCOV}\xspace}
\newcommand{\kasan}{\texttt{KASAN}\xspace}
\newcommand{\syzkaller}{\texttt{syzkaller}\xspace}
\newcommand{\drgn}{\texttt{drgn}\xspace}
\newcommand{\virtme}{\texttt{virtme-ng}\xspace}

% Data-flow tuple notation
\newcommand{\dftuple}[4]{\ensuremath{\langle #1, #2, #3, #4 \rangle}}
\newcommand{\callerid}{\ensuremath{\mathit{PC}_{\mathrm{caller}}}}
\newcommand{\calleeid}{\ensuremath{\mathit{PC}_{\mathrm{callee}}}}

% Listing style
\usepackage{listings}
\lstset{
  basicstyle=\ttfamily\footnotesize,
  breaklines=true,
  frame=single,
  numbers=left,
  numberstyle=\tiny,
  captionpos=b,
}
 for NDSS,
%                       set \anonymousfalse before % shared/macros.tex — Identity-masked macros for both arXiv and NDSS versions
% Toggle anonymization: set \anonymoustrue before % shared/macros.tex — Identity-masked macros for both arXiv and NDSS versions
% Toggle anonymization: set \anonymoustrue before \input{macros} for NDSS,
%                       set \anonymousfalse before \input{macros} for arXiv

\ifanonymous
  \newcommand{\toolname}{{\sc BoundaryFlow}\xspace}
  \newcommand{\anonrepo}{\url{https://anonymous.4open.science/r/XXXX}}
  \newcommand{\authorblock}{%
    \author{Anonymous Submission \\ Paper ID: TBD}%
  }
\else
  \newcommand{\toolname}{{\sc kcov-dataflow}\xspace}
  \newcommand{\anonrepo}{\url{https://github.com/yskzalloc/kcov-dataflow}}
  \newcommand{\authorblock}{%
    \author{Yunseong Kim \\
      Ericsson Software Technology \\
      \texttt{yunseong.kim@est.tech}}%
    % \author{Anonymous Submission}%
  }
\fi

% Technical shorthands
\newcommand{\sancov}{\texttt{SanitizerCoverage}\xspace}
\newcommand{\tracepc}{\texttt{trace-pc}\xspace}
\newcommand{\traceargs}{\texttt{trace-args}\xspace}
\newcommand{\traceret}{\texttt{trace-ret}\xspace}
\newcommand{\kcov}{\texttt{KCOV}\xspace}
\newcommand{\kasan}{\texttt{KASAN}\xspace}
\newcommand{\syzkaller}{\texttt{syzkaller}\xspace}
\newcommand{\drgn}{\texttt{drgn}\xspace}
\newcommand{\virtme}{\texttt{virtme-ng}\xspace}

% Data-flow tuple notation
\newcommand{\dftuple}[4]{\ensuremath{\langle #1, #2, #3, #4 \rangle}}
\newcommand{\callerid}{\ensuremath{\mathit{PC}_{\mathrm{caller}}}}
\newcommand{\calleeid}{\ensuremath{\mathit{PC}_{\mathrm{callee}}}}

% Listing style
\usepackage{listings}
\lstset{
  basicstyle=\ttfamily\footnotesize,
  breaklines=true,
  frame=single,
  numbers=left,
  numberstyle=\tiny,
  captionpos=b,
}
 for NDSS,
%                       set \anonymousfalse before % shared/macros.tex — Identity-masked macros for both arXiv and NDSS versions
% Toggle anonymization: set \anonymoustrue before \input{macros} for NDSS,
%                       set \anonymousfalse before \input{macros} for arXiv

\ifanonymous
  \newcommand{\toolname}{{\sc BoundaryFlow}\xspace}
  \newcommand{\anonrepo}{\url{https://anonymous.4open.science/r/XXXX}}
  \newcommand{\authorblock}{%
    \author{Anonymous Submission \\ Paper ID: TBD}%
  }
\else
  \newcommand{\toolname}{{\sc kcov-dataflow}\xspace}
  \newcommand{\anonrepo}{\url{https://github.com/yskzalloc/kcov-dataflow}}
  \newcommand{\authorblock}{%
    \author{Yunseong Kim \\
      Ericsson Software Technology \\
      \texttt{yunseong.kim@est.tech}}%
    % \author{Anonymous Submission}%
  }
\fi

% Technical shorthands
\newcommand{\sancov}{\texttt{SanitizerCoverage}\xspace}
\newcommand{\tracepc}{\texttt{trace-pc}\xspace}
\newcommand{\traceargs}{\texttt{trace-args}\xspace}
\newcommand{\traceret}{\texttt{trace-ret}\xspace}
\newcommand{\kcov}{\texttt{KCOV}\xspace}
\newcommand{\kasan}{\texttt{KASAN}\xspace}
\newcommand{\syzkaller}{\texttt{syzkaller}\xspace}
\newcommand{\drgn}{\texttt{drgn}\xspace}
\newcommand{\virtme}{\texttt{virtme-ng}\xspace}

% Data-flow tuple notation
\newcommand{\dftuple}[4]{\ensuremath{\langle #1, #2, #3, #4 \rangle}}
\newcommand{\callerid}{\ensuremath{\mathit{PC}_{\mathrm{caller}}}}
\newcommand{\calleeid}{\ensuremath{\mathit{PC}_{\mathrm{callee}}}}

% Listing style
\usepackage{listings}
\lstset{
  basicstyle=\ttfamily\footnotesize,
  breaklines=true,
  frame=single,
  numbers=left,
  numberstyle=\tiny,
  captionpos=b,
}
 for arXiv

\ifanonymous
  \newcommand{\toolname}{{\sc BoundaryFlow}\xspace}
  \newcommand{\anonrepo}{\url{https://anonymous.4open.science/r/XXXX}}
  \newcommand{\authorblock}{%
    \author{Anonymous Submission \\ Paper ID: TBD}%
  }
\else
  \newcommand{\toolname}{{\sc kcov-dataflow}\xspace}
  \newcommand{\anonrepo}{\url{https://github.com/yskzalloc/kcov-dataflow}}
  \newcommand{\authorblock}{%
    \author{Yunseong Kim \\
      Ericsson Software Technology \\
      \texttt{yunseong.kim@est.tech}}%
    % \author{Anonymous Submission}%
  }
\fi

% Technical shorthands
\newcommand{\sancov}{\texttt{SanitizerCoverage}\xspace}
\newcommand{\tracepc}{\texttt{trace-pc}\xspace}
\newcommand{\traceargs}{\texttt{trace-args}\xspace}
\newcommand{\traceret}{\texttt{trace-ret}\xspace}
\newcommand{\kcov}{\texttt{KCOV}\xspace}
\newcommand{\kasan}{\texttt{KASAN}\xspace}
\newcommand{\syzkaller}{\texttt{syzkaller}\xspace}
\newcommand{\drgn}{\texttt{drgn}\xspace}
\newcommand{\virtme}{\texttt{virtme-ng}\xspace}

% Data-flow tuple notation
\newcommand{\dftuple}[4]{\ensuremath{\langle #1, #2, #3, #4 \rangle}}
\newcommand{\callerid}{\ensuremath{\mathit{PC}_{\mathrm{caller}}}}
\newcommand{\calleeid}{\ensuremath{\mathit{PC}_{\mathrm{callee}}}}

% Listing style
\usepackage{listings}
\lstset{
  basicstyle=\ttfamily\footnotesize,
  breaklines=true,
  frame=single,
  numbers=left,
  numberstyle=\tiny,
  captionpos=b,
}
 for arXiv

\ifanonymous
  \newcommand{\toolname}{{\sc BoundaryFlow}\xspace}
  \newcommand{\anonrepo}{\url{https://anonymous.4open.science/r/XXXX}}
  \newcommand{\authorblock}{%
    \author{Anonymous Submission \\ Paper ID: TBD}%
  }
\else
  \newcommand{\toolname}{{\sc kcov-dataflow}\xspace}
  \newcommand{\anonrepo}{\url{https://github.com/yskzalloc/kcov-dataflow}}
  \newcommand{\authorblock}{%
    \author{Yunseong Kim \\
      Ericsson Software Technology \\
      \texttt{yunseong.kim@est.tech}}%
    % \author{Anonymous Submission}%
  }
\fi

% Technical shorthands
\newcommand{\sancov}{\texttt{SanitizerCoverage}\xspace}
\newcommand{\tracepc}{\texttt{trace-pc}\xspace}
\newcommand{\traceargs}{\texttt{trace-args}\xspace}
\newcommand{\traceret}{\texttt{trace-ret}\xspace}
\newcommand{\kcov}{\texttt{KCOV}\xspace}
\newcommand{\kasan}{\texttt{KASAN}\xspace}
\newcommand{\syzkaller}{\texttt{syzkaller}\xspace}
\newcommand{\drgn}{\texttt{drgn}\xspace}
\newcommand{\virtme}{\texttt{virtme-ng}\xspace}

% Data-flow tuple notation
\newcommand{\dftuple}[4]{\ensuremath{\langle #1, #2, #3, #4 \rangle}}
\newcommand{\callerid}{\ensuremath{\mathit{PC}_{\mathrm{caller}}}}
\newcommand{\calleeid}{\ensuremath{\mathit{PC}_{\mathrm{callee}}}}

% Listing style
\usepackage{listings}
\lstset{
  basicstyle=\ttfamily\footnotesize,
  breaklines=true,
  frame=single,
  numbers=left,
  numberstyle=\tiny,
  captionpos=b,
}


\title{Automated Boundary Instrumentation for Deep State-Space Exploration in Operating Systems}
\authorblock
\date{}

\begin{document}
\maketitle

\begin{abstract}
Operating system kernels contain deep state spaces where security-critical
behavior depends on runtime argument values, not merely control-flow topology.
Existing automated exploration tools use control-flow coverage as their sole
feedback signal, creating a fundamental blind spot for value-dependent state
transitions---two executions traversing identical edges may differ only in a
size parameter that determines whether a heap overflow occurs. We present
\toolname{}, a compiler-assisted boundary instrumentation framework that
automatically captures function-level data-flow context---arguments, return
values, and struct field decomposition---with negligible runtime overhead. An
LLVM pass emits lightweight callbacks at function boundaries using debug
metadata for zero-annotation struct expansion; a kernel-resident lock-free
transport delivers records to user-space consumers without interfering with
existing testing infrastructure. We demonstrate that boundary data-flow context
enables both deeper automated state-space exploration and deterministic post-hoc
security analysis. A cross-language pipeline extends instrumentation to Rust
kernel modules by rebuilding the Rust compiler against our modified LLVM
backend---providing native support without post-processing hacks. Evaluated on
five distinct vulnerability patterns across C and Rust kernel modules,
\toolname{} captures information unavailable through any combination of
sanitizer reports, coverage traces, and post-mortem debuggers, with less than
3\% overhead on instrumented paths.
\end{abstract}

\section{Introduction}\label{sec:intro}

Modern operating system kernels form the trusted computing base for all
user-space security guarantees. A single memory corruption in kernel space
compromises process isolation, access control, and cryptographic protections
simultaneously. Automated tools---fuzzers, static analyzers, and
sanitizers---are the primary pre-deployment defense for discovering these
vulnerabilities before adversaries exploit them.

Coverage-guided fuzzing has proven remarkably effective, with tools such as
syzkaller~\cite{vyukov2015syzkaller} discovering thousands of kernel bugs
annually. Yet coverage \emph{saturates}: after initial exploration, new edges
become rare while the remaining bugs hide behind value-dependent state
transitions that produce identical control-flow traces. In
\texttt{io\_uring}'s submission queue processing, the same edges execute for
different opcode values, but only specific opcode sequences trigger
vulnerabilities. In \texttt{binder}'s transaction handling, the
security-relevant difference between a benign and exploitable path lies in
struct field values invisible to edge coverage.

Dynamic taint analysis (e.g., DataFlowSanitizer) tracks byte-level propagation
but imposes 10--100$\times$ overhead---unusable for kernel fuzzing where
millions of executions per second are required. Static taint over-approximates,
producing false positives that drown analysts. What is needed is a lightweight,
always-on signal that captures \emph{what values} flow across function
boundaries without tracking \emph{where every byte propagates}.

Our key insight is that function boundaries are natural chokepoints in kernel
execution. Arguments encode the \emph{intent} of a call; return values encode
the \emph{outcome}. Capturing these at compile time---with zero source
annotation---via a lock-free kernel transport gives both fuzzers and analysts a
new dimension of observability at negligible cost.

The Linux kernel now integrates Rust as a second systems language, adding
cross-language FFI boundaries where traditional debugging tools fail entirely.
The Rust compiler (\texttt{rustc}) at optimization level \texttt{-O2} elides
DWARF variable location information, making post-mortem debuggers like
\texttt{drgn} unable to read function arguments from vmcore dumps. Our
framework addresses this by instrumenting at the LLVM IR level---below the
language frontend but above the optimizer---capturing arguments before they are
optimized away.

We make the following contributions:
\begin{itemize}[nosep]
  \item A zero-annotation compiler instrumentation that extracts structured
        data-flow context at function boundaries, including automatic
        decomposition of aggregate types via DWARF debug metadata
        (\texttt{DICompositeType}), requiring no source modification.
  \item A kernel-resident lock-free transport mechanism delivering data-flow
        records to user-space with zero interference to existing kernel testing
        infrastructure (KCOV, syzkaller), safe in process, softirq, and hardirq
        contexts.
  \item A cross-language instrumentation approach: rebuilding \texttt{rustc}
        against our modified LLVM backend gives Rust kernel modules native
        dataflow instrumentation support---no post-compilation pipeline needed.
  \item Empirical demonstration on five distinct vulnerability patterns (OOB,
        UAF, double-free, 10-deep call chain, Rust FFI) showing that boundary
        data-flow context provides information unavailable through any
        combination of KASAN, drgn, and edge coverage.
  \item A complete end-to-end verification methodology using \texttt{virtme-ng},
        \texttt{drgn}, and KASAN to validate correctness of captured data
        against ground truth.
\end{itemize}


\section{Threat Model \& Problem Statement}\label{sec:threat}

\subsection{Threat Model}

We consider a standard kernel fuzzing adversary: an unprivileged local user
with access to the syscall interface who attempts to trigger memory safety
violations in kernel subsystems. The defense goal is to maximize the
probability of discovering such vulnerabilities during pre-deployment testing.

We assume the kernel is compiled with sanitizer support (KASAN) and debug
information (\texttt{-g}), which is standard practice for fuzzing builds. The
attacker does not have physical access or the ability to modify kernel code.

\subsection{Problem Statement}

Let $\mathcal{K}$ be a kernel with state space $S = \{s_1, \ldots, s_n\}$ and
transition function $T: S \times I \to S$ where input $I$ triggers state
changes via syscalls. Coverage-guided fuzzing explores states reachable via
distinct control-flow paths $P \subset S$. However, $|P| \ll |S|$ because many
distinct states share identical control-flow paths---they differ only in
argument values passed between functions.

Define the \emph{boundary context} $B(f, i) = (\mathit{args}(f, i),
\mathit{ret}(f, i))$ for function $f$ on input $i$. Our goal: expose $B$ to
the exploration engine so it can distinguish states in $S \setminus P$ that are
invisible to pure control-flow coverage.

\paragraph{Motivating Example.} Consider a kernel module with function
\texttt{vuln\_process(struct session\_data *sd, unsigned int size)}. When
\texttt{size=15}, the function operates within bounds. When \texttt{size=32},
it overflows a 16-byte buffer. Both executions traverse \emph{identical} basic
blocks---the same \texttt{memset()} call, the same return path. Edge coverage
cannot distinguish them. But boundary instrumentation captures
\texttt{size=15} vs.\ \texttt{size=32} at function entry, immediately revealing
the value that triggers the overflow.

\subsection{Requirements}

\begin{description}[nosep]
  \item[R1] Zero source annotation (must work on unmodified kernel source).
  \item[R2] $<$5\% overhead on instrumented paths.
  \item[R3] No interference with existing KCOV/syzkaller infrastructure.
  \item[R4] Safe in interrupt context (process, softirq, hardirq).
  \item[R5] Cross-language support (C and Rust kernel modules).
  \item[R6] Struct field decomposition (capture individual fields, not just
        the pointer).
\end{description}

\section{Design}\label{sec:design}

\subsection{Architectural Overview}

\toolname{} operates as a three-layer pipeline (Figure~\ref{fig:arch}):

\begin{enumerate}[nosep]
  \item \textbf{Compile-time instrumentation:} An LLVM pass inserts data-flow
        callbacks at function boundaries, with struct field expansion driven by
        DWARF debug metadata.
  \item \textbf{Runtime transport:} A lock-free per-task ring buffer in kernel
        space receives records from the callbacks.
  \item \textbf{User-space consumer:} A fuzzer executor or analyst tool reads
        the buffer via \texttt{mmap()} for zero-copy access.
\end{enumerate}

\begin{figure}[t]
\centering
\fbox{\parbox{0.9\columnwidth}{\small
\textbf{Layer 1: Compile-Time} \\
\texttt{clang -fsanitize-coverage=dataflow-args,dataflow-ret} \\
$\downarrow$ inserts callbacks at function entry/exit \\[0.5em]
\textbf{Layer 2: Kernel Runtime} \\
\texttt{\_\_sanitizer\_cov\_trace\_args(pc, idx, sz, ptr, offsets, nf)} \\
$\downarrow$ atomic write to per-task ring buffer \\[0.5em]
\textbf{Layer 3: User-Space} \\
\texttt{open(/dev/kcov\_dataflow)} $\to$ \texttt{ioctl(INIT)} $\to$ \texttt{mmap()} $\to$ read
}}
\caption{Three-layer architecture of \toolname{}.}\label{fig:arch}
\end{figure}

The system is activated entirely by compiler flags---no source annotation,
no kernel configuration beyond enabling the Kconfig options
\texttt{CONFIG\_KCOV\_DATAFLOW\_ARGS} and \texttt{CONFIG\_KCOV\_DATAFLOW\_RET}.

\subsection{Compile-Time Instrumentation}\label{sec:compile}

The instrumentation pass resides in LLVM's existing \sancov framework
(\texttt{llvm/lib/Transforms/Instrumentation/SanitizerCoverage.cpp}), activated
by the flags \texttt{-fsanitize-coverage=dataflow-args,dataflow-ret}.

\subsubsection{Argument Capture (InjectTraceForArgs)}

For each function $F$ with a \texttt{DISubprogram} debug info attachment, the
pass iterates over formal parameters:

\begin{enumerate}[nosep]
  \item \textbf{Type Resolution:} Resolve each parameter's
        \texttt{DILocalVariable} to its \texttt{DIType}. Strip typedefs and
        qualifiers via a recursive \texttt{stripDITypedefs()} helper.
  \item \textbf{Composite Detection:} If the resolved type is a
        \texttt{DICompositeType} (struct, union, or class), walk its member
        list (\texttt{DIDerivedType} elements) to extract:
        \begin{itemize}[nosep]
          \item Field byte-offset within the struct
          \item Field size in bytes
        \end{itemize}
  \item \textbf{Type Hashing:} Compute an FNV-1a hash of the composite type's
        name. This hash is stored at index~0 of the offset array, enabling
        user-space consumers to deduplicate struct types across compilation
        units without string comparison.
  \item \textbf{Offset Array Emission:} Create a module-level
        \texttt{ConstantArray} global:
        \texttt{@\_\_sancov\_offsets\_ = [hash, off$_0$, sz$_0$, off$_1$,
        sz$_1$, \ldots, off$_N$, sz$_N$]}
  \item \textbf{Callback Insertion:} At function entry, emit:
\begin{lstlisting}[language=C,numbers=none]
__sanitizer_cov_trace_args(
  pc,          // function address
  arg_idx,     // parameter index (0-based)
  arg_size,    // sizeof(parameter)
  ptr,         // pointer to argument value
  offsets_ptr, // &offsets[1] (past hash)
  num_fields   // number of struct fields
);
\end{lstlisting}
\end{enumerate}

For scalar arguments (integers, pointers without struct type), the pass spills
the value to a stack \texttt{alloca} so the kernel callback receives a uniform
pointer interface. \texttt{offsets\_ptr} is \texttt{null} and
\texttt{num\_fields} is~0 for scalars.

\subsubsection{Return Value Capture (InjectTraceForRet)}

Using LLVM's \texttt{EscapeEnumerator}, the pass locates every
\texttt{ReturnInst} in $F$ and inserts:
\begin{lstlisting}[language=C,numbers=none]
__sanitizer_cov_trace_ret(
  pc, ret_size, ret_ptr, offsets_ptr, num_fields
);
\end{lstlisting}

This captures the return value \emph{after} the function body executes but
\emph{before} the caller receives it---the ideal observation point for
detecting corruption that occurred during execution.

\subsubsection{Edge Cases and Graceful Degradation}

\begin{itemize}[nosep]
  \item \textbf{No debug info:} Function skipped entirely (no crash, no
        partial instrumentation).
  \item \textbf{Variadic functions:} Skipped (va\_list arguments have no
        static type information).
  \item \textbf{Naked functions:} Skipped (no prologue/epilogue to instrument).
  \item \textbf{Indirect calls:} The callee is instrumented at its own entry
        point regardless of how it was called.
  \item \textbf{Inlined functions:} LLVM's inliner runs before our pass in the
        default pipeline; inlined functions are not separately instrumented
        (their arguments become part of the caller's body).
\end{itemize}

\subsection{Runtime Transport Layer}\label{sec:transport}

\subsubsection{Per-Task State}

The kernel backend adds three fields to \texttt{struct task\_struct}:
\begin{lstlisting}[language=C,numbers=none]
u64 *kcov_df_area;    // mmap'd ring buffer
unsigned int kcov_df_size;  // capacity (words)
atomic_t kcov_df_seq; // monotonic counter
\end{lstlisting}

\subsubsection{Write Path (Critical Section)}

The callback implementation follows a strict safety protocol:

\begin{algorithmic}[1]
\Function{\_\_sanitizer\_cov\_trace\_args}{pc, idx, sz, ptr, offs, nf}
  \If{$\neg$current$\to$kcov\_df\_area \textbf{or} in\_nmi()} \Return \EndIf
  \State $\mathit{needed} \gets 6 + \mathit{nf}$ \Comment{record size in u64 words}
  \State $\mathit{slot} \gets$ atomic\_add(needed, \&seq) \Comment{reserve slot}
  \If{$\mathit{slot} + \mathit{needed} > \mathit{size}$} \Return \EndIf \Comment{drop if full}
  \State area[slot+0] $\gets$ seq\_nr
  \State area[slot+1] $\gets$ ENTRY\_FLAG
  \State area[slot+2] $\gets$ pc
  \State area[slot+3] $\gets$ (idx $\ll$ 32) | sz
  \State area[slot+4] $\gets$ nf
  \For{$i = 0$ \textbf{to} $\mathit{nf}-1$}
    \State $\mathit{field\_ptr} \gets$ ptr + offs[$2i$] \Comment{byte offset}
    \State copy\_from\_kernel\_nofault(\&area[slot+5+i], field\_ptr, offs[$2i+1$])
  \EndFor
\EndFunction
\end{algorithmic}

Key safety properties:
\begin{itemize}[nosep]
  \item \textbf{No locks:} \texttt{atomic\_add} is a single x86
        \texttt{LOCK XADD} instruction.
  \item \textbf{No allocation:} Buffer pre-allocated at \texttt{ioctl(INIT)}.
  \item \textbf{No blocking:} Overflow silently drops events.
  \item \textbf{Fault tolerance:} \texttt{copy\_from\_kernel\_nofault()} safely
        handles unmapped pages, KASAN red zones, and freed slab objects.
  \item \textbf{ERR\_PTR handling:} Pointer values in $[-4095, -1]$ are
        detected and skipped (kernel error codes, not valid addresses).
  \item \textbf{Recursion prevention:} Callbacks marked \texttt{notrace},
        \texttt{\_\_no\_sanitize\_coverage}, \texttt{noinline}.
  \item \textbf{No printk:} Zero I/O in the data path.
\end{itemize}

\subsubsection{Device Interface}

\texttt{/sys/kernel/debug/kcov\_dataflow} is \emph{completely independent} from
the legacy \texttt{/sys/kernel/debug/kcov} device. This is a critical design
decision: syzkaller's executor opens and manages KCOV; our device must not
interfere.

\begin{table}[t]
\centering
\caption{Ioctl interface specification.}\label{tab:ioctl}
\small
\begin{tabular}{@{}lll@{}}
\toprule
\textbf{Ioctl} & \textbf{Value} & \textbf{Action} \\
\midrule
\texttt{KCOV\_DF\_INIT\_TRACE} & \texttt{0x80086401} & Allocate buffer \\
\texttt{KCOV\_DF\_ENABLE} & \texttt{0x6464} & Start recording \\
\texttt{KCOV\_DF\_DISABLE} & \texttt{0x6465} & Stop recording \\
\bottomrule
\end{tabular}
\end{table}

The ioctl namespace uses magic \texttt{'d'} (Table~\ref{tab:ioctl}). User-space
maps the buffer via \texttt{mmap()} for zero-copy reads---no syscall overhead
to retrieve records.

\subsection{Cross-Language Support: Rust}\label{sec:rust}

\subsubsection{The Rust LLVM Problem}

The Rust compiler (\texttt{rustc}) bundles its own LLVM library (version~22 for
rustc~1.95). Our custom \texttt{cl::opt} flags
(\texttt{-sanitizer-coverage-dataflow-args/ret}) do not exist in this bundled
LLVM. Passing them via \texttt{-Cllvm-args} results in:
\begin{lstlisting}[numbers=none]
Unknown command line argument
  '-sanitizer-coverage-dataflow-args'
\end{lstlisting}

\subsubsection{Solution: Rebuild rustc with Our LLVM}

The Rust build system (\texttt{x.py}) supports linking against an external LLVM
via \texttt{config.toml}:
\begin{lstlisting}[language=bash,numbers=none]
[llvm]
download-ci-llvm = false

[target.x86_64-unknown-linux-gnu]
llvm-config = "/path/to/our/llvm-config"
\end{lstlisting}

Building stage~1 (\texttt{python3 x.py build --stage 1 library}) produces a
\texttt{rustc} binary linked against our LLVM~23 with the dataflow pass. This
rustc natively supports:
\begin{lstlisting}[language=bash,numbers=none]
rustc -Cpasses=sancov-module \
  -Cllvm-args="-sanitizer-coverage-level=3 \
    -sanitizer-coverage-dataflow-args \
    -sanitizer-coverage-dataflow-ret" \
  --emit=obj module.rs
\end{lstlisting}

The resulting object file contains references to
\texttt{\_\_sanitizer\_cov\_trace\_args} and
\texttt{\_\_sanitizer\_cov\_trace\_ret}---identical to C modules compiled with
our clang.

\subsubsection{Kernel Build Integration}

We add \texttt{RUSTFLAGS\_KCOV\_DATAFLOW} to the kernel's
\texttt{scripts/Makefile.kcov}:
\begin{lstlisting}[language=make,numbers=none]
export RUSTFLAGS_KCOV_DATAFLOW := \
  -Cpasses=sancov-module \
  -Cllvm-args=-sanitizer-coverage-level=3 \
  -Cllvm-args=-sanitizer-coverage-dataflow-args \
  -Cllvm-args=-sanitizer-coverage-dataflow-ret \
  -Cdebuginfo=2
\end{lstlisting}

\texttt{scripts/Makefile.lib} applies these flags when a module declares
\texttt{KCOV\_DATAFLOW\_module.o := y}---the same single-line opt-in as C
modules.

\subsubsection{Alternative: Post-Compilation Pipeline}

For environments where rebuilding rustc is impractical, we provide a fallback:
\begin{enumerate}[nosep]
  \item \texttt{rustc --emit=llvm-ir -g}: emit LLVM IR text with debug info
  \item \texttt{opt -passes=sancov-module ...}: instrument with our pass
  \item \texttt{llc -filetype=obj}: compile to object
\end{enumerate}
The text format (\texttt{.ll}) avoids bitcode version incompatibilities between
LLVM~22 (rustc's) and LLVM~23 (ours).


\section{Implementation}\label{sec:impl}

\subsection{LLVM Pass Implementation (600~LOC)}

The pass extends \texttt{SanitizerCoverage.cpp} with two primary methods and
two helpers. We detail the integration points across the LLVM/Clang codebase.

\subsubsection{Driver Integration}

Enabling the pass requires changes to seven files:

\begin{table}[t]
\centering
\caption{Files modified in LLVM/Clang for dataflow coverage.}\label{tab:files}
\small
\begin{tabular}{@{}lp{5cm}@{}}
\toprule
\textbf{File} & \textbf{Change} \\
\midrule
\texttt{CodeGenOptions.def} & Add \texttt{CODEGENOPT} flags \\
\texttt{CodeGenOptions.h} & Include in \texttt{hasSanitizeCoverage()} \\
\texttt{Options.td} & Define \texttt{-cc1} flags \\
\texttt{SanitizerArgs.cpp} & Parse \texttt{"dataflow-args"}/\texttt{"dataflow-ret"}, enum values $1\ll20$, $1\ll21$ \\
\texttt{BackendUtil.cpp} & Wire to \texttt{SanitizerCoverageOptions} \\
\texttt{Instrumentation.h} & Add \texttt{DataflowArgs}/\texttt{DataflowRet} bools \\
\texttt{SanitizerCoverage.cpp} & Main implementation \\
\bottomrule
\end{tabular}
\end{table}

The flags imply edge coverage (\texttt{-fsanitize-coverage=edge}) to ensure
basic block instrumentation is also present, enabling correlation between
data-flow events and control-flow traces.

\subsubsection{Helper: stripDITypedefs}

\begin{lstlisting}[language=C++,caption={Recursive typedef/qualifier stripping.}]
static DIType *stripDITypedefs(DIType *T) {
  while (T) {
    if (auto *DT = dyn_cast<DIDerivedType>(T)) {
      unsigned Tag = DT->getTag();
      if (Tag == dwarf::DW_TAG_typedef ||
          Tag == dwarf::DW_TAG_const_type ||
          Tag == dwarf::DW_TAG_volatile_type ||
          Tag == dwarf::DW_TAG_restrict_type)
        T = DT->getBaseType();
      else break;
    } else break;
  }
  return T;
}
\end{lstlisting}

\subsubsection{Helper: getStructFieldOffsets}

This function extracts the memory layout of a composite type:
\begin{enumerate}[nosep]
  \item Iterate \texttt{DICompositeType::getElements()}
  \item For each \texttt{DIDerivedType} member with tag
        \texttt{DW\_TAG\_member}: extract \texttt{getOffsetInBits()/8} and
        \texttt{getSizeInBits()/8}
  \item Compute FNV-1a hash: \texttt{hash = 0xcbf29ce484222325; for each byte
        of name: hash \^{}= byte; hash *= 0x100000001b3}
  \item Return array: \texttt{[hash, off$_0$, sz$_0$, \ldots]}
\end{enumerate}

\subsubsection{InjectTraceForArgs}

\begin{algorithmic}[1]
\Function{InjectTraceForArgs}{$F$}
  \State $\mathit{SP} \gets F.\mathrm{getSubprogram()}$
  \If{$\mathit{SP} = \mathrm{null}$} \Return \EndIf
  \State $\mathit{EntryBB} \gets F.\mathrm{getEntryBlock()}$
  \State $\mathit{IP} \gets$ first non-alloca instruction in EntryBB
  \For{$i = 0$ \textbf{to} $F.\mathrm{arg\_size()}-1$}
    \State $\mathit{Arg} \gets F.\mathrm{getArg}(i)$
    \State $T \gets \mathrm{resolveArgType}(\mathit{SP}, i)$
    \State $T \gets \mathrm{stripDITypedefs}(T)$
    \If{$T$ is \texttt{DICompositeType}}
      \State $(\mathit{Offsets}, \mathit{NF}) \gets \mathrm{getStructFieldOffsets}(T)$
      \State $G \gets$ create global constant array from Offsets
      \State $\mathit{OffsetsPtr} \gets$ GEP($G$, 0, 1) \Comment{skip hash}
    \Else
      \State $\mathit{OffsetsPtr} \gets \mathrm{null}$; $\mathit{NF} \gets 0$
    \EndIf
    \If{Arg is pointer type}
      \State $\mathit{Ptr} \gets \mathrm{Arg}$
    \Else
      \State $\mathit{Alloca} \gets$ CreateAlloca(Arg.getType())
      \State Store Arg $\to$ Alloca
      \State $\mathit{Ptr} \gets$ Alloca
    \EndIf
    \State Insert call: trace\_args(PC, $i$, sizeof(Arg), Ptr, OffsetsPtr, NF)
  \EndFor
\EndFunction
\end{algorithmic}

\subsubsection{InjectTraceForRet}

Uses \texttt{EscapeEnumerator} to find all function exit points (including
those created by exception handling). For each \texttt{ReturnInst}:
\begin{enumerate}[nosep]
  \item Extract the return value
  \item If void return: skip
  \item If struct return: resolve type, get offsets (same as args)
  \item Spill return value to alloca if scalar
  \item Insert \texttt{\_\_sanitizer\_cov\_trace\_ret(pc, ret\_size, ptr,
        offsets, nf)}
\end{enumerate}

\subsection{Kernel Backend (400~LOC)}\label{sec:kernel-impl}

\subsubsection{Data Structures}

\begin{lstlisting}[language=C,caption={Kernel-side data structures (simplified).}]
struct kcov_dataflow {
  u64 *area;          // vmalloc'd buffer
  unsigned int size;  // capacity in u64 words
  atomic_t seq;       // monotonic sequence
  enum { DF_DISABLED, DF_ENABLED } mode;
  struct task_struct *task; // owning task
};
\end{lstlisting}

\subsubsection{Callback Implementation}

\begin{lstlisting}[language=C,caption={Kernel trace\_args callback (simplified).}]
void notrace __no_sanitize_coverage noinline
__sanitizer_cov_trace_args(u64 pc, u32 arg_idx,
    u32 arg_size, void *ptr,
    u64 *offsets, u32 num_fields)
{
  struct task_struct *t = current;
  u64 *area = t->kcov_df_area;
  unsigned int size = t->kcov_df_size;
  unsigned int needed, pos;
  u64 val;
  int i;

  if (unlikely(!area) || in_nmi())
    return;

  needed = 6 + num_fields;
  pos = atomic_add_return(needed,
          &t->kcov_df_seq) - needed;
  if (pos + needed > size)
    return; // buffer full, drop

  area[pos+0] = atomic_read(&t->kcov_df_seq);
  area[pos+1] = 1; // ENTRY flag
  area[pos+2] = pc;
  area[pos+3] = ((u64)arg_idx << 32) | arg_size;
  area[pos+4] = num_fields;

  if (num_fields && offsets && ptr) {
    for (i = 0; i < num_fields; i++) {
      void *field = ptr + offsets[2*i];
      u32 fsz = (u32)offsets[2*i + 1];
      val = 0;
      if (!IS_ERR_VALUE((unsigned long)field))
        copy_from_kernel_nofault(&val,
          field, min(fsz, 8u));
      area[pos+5+i] = val;
    }
  } else if (ptr) {
    val = 0;
    if (!IS_ERR_VALUE((unsigned long)ptr))
      copy_from_kernel_nofault(&val, ptr,
        min(arg_size, 8u));
    area[pos+5] = val;
  }
}
EXPORT_SYMBOL(__sanitizer_cov_trace_args);
\end{lstlisting}

\subsubsection{File Operations}

The device registers its own \texttt{file\_operations} with \texttt{open},
\texttt{release}, \texttt{unlocked\_ioctl}, and \texttt{mmap} handlers. The
\texttt{mmap} handler maps the \texttt{vmalloc}'d buffer directly into
user-space address space.

\subsection{Build System Integration}

\subsubsection{C Modules}

\texttt{scripts/Makefile.kcov} exports:
\begin{lstlisting}[language=make,numbers=none]
kcov-dataflow-flags-y := \
  -fsanitize-coverage=dataflow-args,dataflow-ret \
  -g -fno-inline
export CFLAGS_KCOV_DATAFLOW := \
  $(kcov-dataflow-flags-y)
\end{lstlisting}

\texttt{scripts/Makefile.lib} applies these when the module Makefile declares:
\begin{lstlisting}[language=make,numbers=none]
KCOV_DATAFLOW_my_module.o := y
\end{lstlisting}

\subsubsection{Rust Modules}

\texttt{scripts/Makefile.kcov} additionally exports:
\begin{lstlisting}[language=make,numbers=none]
export RUSTFLAGS_KCOV_DATAFLOW := \
  -Cpasses=sancov-module \
  -Cllvm-args=-sanitizer-coverage-level=3 \
  -Cllvm-args=-sanitizer-coverage-dataflow-args \
  -Cllvm-args=-sanitizer-coverage-dataflow-ret \
  -Cdebuginfo=2
\end{lstlisting}

\texttt{scripts/Makefile.lib} applies \texttt{RUSTFLAGS\_KCOV\_DATAFLOW} to
\texttt{\_rust\_flags} using the same \texttt{KCOV\_DATAFLOW\_module.o}
mechanism---unified opt-in for both languages.

\subsection{User-Space Consumer}

A minimal consumer program (our ``trigger'' tool):
\begin{lstlisting}[language=C,caption={User-space kcov\_dataflow consumer.}]
int fd = open("/dev/kcov_dataflow", O_RDWR);
ioctl(fd, KCOV_DF_INIT_TRACE, BUF_SIZE);
u64 *buf = mmap(NULL, BUF_SIZE * 8,
  PROT_READ, MAP_SHARED, fd, 0);
ioctl(fd, KCOV_DF_ENABLE, 0);

// Trigger kernel code path
write(trigger_fd, "x", 1);

ioctl(fd, KCOV_DF_DISABLE, 0);
// Parse records from buf[0..seq]
\end{lstlisting}


\section{Evaluation}\label{sec:eval}

We structure our evaluation around five research questions:
\begin{description}[nosep]
  \item[RQ1] Does boundary instrumentation correctly capture data-flow context?
  \item[RQ2] Does it provide information unavailable through existing tools?
  \item[RQ3] What is the runtime overhead?
  \item[RQ4] Does cross-language (Rust) support work correctly?
  \item[RQ5] Can the captured data enable root-cause analysis?
\end{description}

\subsection{Experimental Setup}\label{sec:setup}

\begin{table}[t]
\centering
\caption{Experimental environment.}\label{tab:env}
\small
\begin{tabular}{@{}ll@{}}
\toprule
\textbf{Component} & \textbf{Version/Configuration} \\
\midrule
Kernel & linux-next 7.1.0-rc4 \\
C Compiler & Custom clang/LLVM 23 (trunk + dataflow pass) \\
Rust Compiler & rustc 1.98.0-nightly (built against our LLVM 23) \\
Kernel Config & KASAN=y, KCOV=y, KCOV\_DATAFLOW\_ARGS=y \\
              & KCOV\_DATAFLOW\_RET=y, RUST=y, DEBUG\_INFO\_DWARF5=y \\
Runtime & virtme-ng (QEMU/KVM), 4 vCPUs, 4\,GB RAM \\
Debugger & drgn 0.2.0 (vmcore + /proc/kcore analysis) \\
\bottomrule
\end{tabular}
\end{table}

We built five purpose-built kernel modules exercising distinct vulnerability
classes (Table~\ref{tab:env}). Each module exposes a \texttt{procfs} or
\texttt{debugfs} trigger that can be activated from user-space while
\texttt{kcov\_dataflow} is recording.

\subsection{RQ1: Correctness Verification}\label{sec:rq1}

\subsubsection{Case 1: Heap Out-of-Bounds Write}

The test module allocates a 16-byte buffer and calls
\texttt{vuln\_process(sd, size)} where \texttt{size} is attacker-controlled.
When \texttt{size=32}, \texttt{memset()} overflows the buffer.

\paragraph{kcov\_dataflow ENTRY capture:}
\begin{lstlisting}[numbers=none]
[ENTRY] seq=5 pc=vuln_process arg[0](8)
  struct session_data:
    .id    = 0x1337
    .buf   = "initial_"
    .size  = 15
\end{lstlisting}

\paragraph{kcov\_dataflow RET capture:}
\begin{lstlisting}[numbers=none]
[RET] seq=6 pc=vuln_process ret(8)
  struct session_data:
    .id    = 0x1337
    .buf   = (corrupted)
    .size  = 32  [OVERFLOW]
\end{lstlisting}

\paragraph{KASAN report:}
\begin{lstlisting}[numbers=none]
BUG: KASAN: slab-out-of-bounds in
  vuln_process+0x43/0xb0 [simple_vuln_mod]
Write of size 32 at addr ffff888001f09d10
\end{lstlisting}

\paragraph{Analysis:} The ENTRY record shows the struct in its valid state
(\texttt{size=15}). The RET record shows \texttt{size=32}---the corrupted
value. KASAN confirms the overflow but only reports the \emph{address}; our
tool identifies \emph{which field} was corrupted and \emph{what value} caused
the overflow.

\subsubsection{Case 2: Use-After-Free Write}

The module frees a slab object, then writes attacker-controlled data to it.

\paragraph{kcov\_dataflow ENTRY:}
\begin{lstlisting}[numbers=none]
[ENTRY] seq=5 pc=uaf_write arg[0](8)
  struct session_data:
    .id    = 0x0  (freed, zero-init)
    .buf   = 0xfdfdfdfd... (KASAN poison)
    .size  = 16
\end{lstlisting}

\paragraph{kcov\_dataflow RET:}
\begin{lstlisting}[numbers=none]
[RET] seq=6 pc=uaf_write ret(8)
  struct session_data:
    .id    = 0x41414141  [ATTACKER VALUE]
    .buf   = "UAF_CORR"
    .size  = 0xdead
\end{lstlisting}

\paragraph{Analysis:} The ENTRY record reveals KASAN poison bytes
(\texttt{0xfd}) in the buffer---proof that the object was freed. The RET record
shows attacker-controlled values written to the freed object. This provides
\emph{controllability evidence} for exploit development: the attacker can write
arbitrary values to specific struct fields.

\subsubsection{Case 3: Double-Free Write}

Similar to UAF but the object is freed twice before being written.

\paragraph{kcov\_dataflow captures:}
\begin{lstlisting}[numbers=none]
ENTRY: {id=0, buf=POISON(0xff), size=99}
RET:   {id=0xdf00df00, buf="DOUBLEFR",
        size=0xdf}
\end{lstlisting}

The \texttt{0xff} poison pattern (distinct from UAF's \texttt{0xfd}) indicates
a double-free state in KASAN's slab allocator.

\subsubsection{Case 4: 10-Deep Call Chain Propagation}

This module implements a 10-function deep call chain where a tainted offset
propagates through multiple transformations:
\begin{lstlisting}[numbers=none]
parse_header() -> validate_header()
  -> lookup_slot() -> transform_offset()
  -> ... (6 more levels)
  -> commit_write(idx)  // OOB if idx >= 8
\end{lstlisting}

\paragraph{kcov\_dataflow trace (abbreviated):}
\begin{lstlisting}[numbers=none]
[ENTRY] seq=1 pc=parse_header arg[1]=16
[ENTRY] seq=3 pc=validate_header arg[0]=48
[ENTRY] seq=5 pc=lookup_slot arg[0]=24
[ENTRY] seq=7 pc=transform_offset arg[0]=8
  ...
[ENTRY] seq=19 pc=commit_write arg[0]=8
[RET]   seq=20 pc=commit_write ret=-EFAULT
\end{lstlisting}

\paragraph{Root-cause analysis:} The trace shows offset propagation:
$16 \to 48 \to 24 \to 8$. The bug is in \texttt{validate\_header()} which
uses \texttt{\%16} instead of \texttt{\%8} for bounds checking, allowing
\texttt{commit\_write()} to receive \texttt{idx=8} (one past the valid range).
This analysis is \emph{impossible} from KASAN alone, which only reports the
final OOB access.

\subsubsection{Case 5: Rust Kernel Module}

The Rust module \texttt{rust\_verify\_mod} implements:
\begin{lstlisting}[numbers=none]
#[repr(C)]
pub struct RustData {
    pub id: u32,
    pub value: u64,
    pub flags: u32,
}

#[no_mangle]
pub extern "C" fn rust_process_data(
    data: *mut RustData,
    multiplier: u32,
    cookie: u64
) -> u64 { ... }
\end{lstlisting}

Called with known values: \texttt{data=\{id=0xcafe, value=0xdeadbeef12345678,
flags=0xa\}}, \texttt{multiplier=3}, \texttt{cookie=0x1111111111111111}.

\paragraph{kcov\_dataflow capture:}
\begin{lstlisting}[numbers=none]
[ENTRY] seq=9 pc=rust_process_data arg[0](8)
  struct RustData:
    .id    = 0xcafe
    .value = 0xdeadbeef12345678
    .flags = 0xa
[ENTRY] seq=10 pc=rust_process_data arg[1](4)
    = 0x3  (multiplier)
[ENTRY] seq=11 pc=rust_process_data arg[2](8)
    = 0x1111111111111111  (cookie)
[RET] seq=75 pc=rust_process_data ret(8)
    = 0xefbed00023456789
\end{lstlisting}

\paragraph{Verification:} \texttt{printk} ground truth confirms:
\texttt{id=0x260fa} ($\texttt{0xcafe} \times 3$),
\texttt{value=0xefbed00023456789} ($\texttt{0xdeadbeef12345678} +
\texttt{0x1111111111111111}$), \texttt{flags=0xfa} ($\texttt{0xa} |
\texttt{0xf0}$). All captured values match exactly.

\subsection{RQ2: Information Gain Over Existing Tools}\label{sec:rq2}

\begin{table*}[t]
\centering
\caption{Detailed observability comparison: what each tool reveals for each vulnerability case.}\label{tab:detailed-gain}
\small
\begin{tabular}{@{}l|ccc|ccc|ccc@{}}
\toprule
& \multicolumn{3}{c|}{\textbf{KASAN}} & \multicolumn{3}{c|}{\textbf{drgn vmcore}} & \multicolumn{3}{c}{\textbf{\toolname{}}} \\
\textbf{Case} & Detect & Which arg & Value & Detect & Which arg & Value & Detect & Which arg & Value \\
\midrule
OOB Write & \cmark & \xmark & \xmark & \cmark & \cmark & \cmark & \cmark & \cmark & \cmark \\
UAF Write & \cmark & \xmark & \xmark & \cmark & \cmark & \cmark & \cmark & \cmark & \cmark \\
Double-Free & \cmark & \xmark & \xmark & \cmark & \cmark & \cmark & \cmark & \cmark & \cmark \\
Deep Chain & \cmark & \xmark & \xmark & partial & manual & manual & \cmark & \cmark & \cmark \\
Rust Module & N/A & N/A & N/A & \xmark & \xmark & \xmark & \cmark & \cmark & \cmark \\
\midrule
Silent Logic & \xmark & \xmark & \xmark & \xmark & \xmark & \xmark & \cmark & \cmark & \cmark \\
\bottomrule
\end{tabular}
\end{table*}

Table~\ref{tab:detailed-gain} shows the detailed comparison. Key findings:

\begin{itemize}[nosep]
  \item \textbf{KASAN} detects memory violations but provides only the
        \emph{address} of the access. It cannot identify which function
        argument carried the bad value, what the value was, or how it
        propagated through the call chain.
  \item \textbf{drgn vmcore} can read C locals from a crash dump, but:
        (a)~requires a crash (misses silent bugs), (b)~requires manual
        navigation of stack frames, (c)~fails entirely on Rust modules at
        \texttt{-O2} due to DWARF variable location elision.
  \item \textbf{\toolname{}} captures all information automatically, works for
        both C and Rust, does not require a crash, and provides temporal
        ordering (ENTRY before corruption, RET after).
\end{itemize}

\paragraph{The Rust Observability Gap.} This is our most significant finding:
for Rust kernel modules compiled at production optimization levels, \toolname{}
is the \emph{only} method that can observe function arguments at runtime.
\texttt{drgn} analysis of a Rust vmcore shows:
\begin{lstlisting}[numbers=none]
>>> prog.stack_trace(task)
#0 rust_process_data+0x380/0x380
>>> frame.locals()
{}  # EMPTY - no variables visible
\end{lstlisting}

This is because \texttt{rustc -O2} does not emit \texttt{DW\_AT\_location}
attributes for function parameters---the DWARF standard allows this when
values are kept only in registers that are subsequently clobbered.

\subsection{RQ3: Performance Overhead}\label{sec:rq3}

We measure overhead using a dedicated micro-benchmark module containing 8
functions with 1--8 arguments each (36 argument captures + 8 return captures =
44 callbacks per iteration). We run 5{,}000 iterations in QEMU/KVM, taking the
best of 3 runs after 3 warmup rounds to eliminate cache effects.

\begin{table}[t]
\centering
\caption{Performance overhead measurement (8-function module, 44 callbacks/iter).}\label{tab:overhead}
\small
\begin{tabular}{@{}lrr@{}}
\toprule
\textbf{Configuration} & \textbf{Per-iter latency} & \textbf{Overhead} \\
\midrule
No recording (baseline) & 15.01\,$\mu$s & --- \\
\toolname{} recording ON & 16.26\,$\mu$s & +8.3\% \\
\bottomrule
\end{tabular}
\end{table}

\paragraph{Per-callback cost breakdown.}
The absolute cost is 1.24\,$\mu$s per iteration for 44 callbacks, yielding
$\sim$\textbf{27\,ns per callback}. This decomposes as:
\begin{itemize}[nosep]
  \item One \texttt{LOCK XADD} instruction ($\sim$20 cycles on x86-64)
  \item One \texttt{copy\_from\_kernel\_nofault()} per struct field ($\sim$5\,ns each)
  \item Branch check \texttt{if (!area)} ($\sim$1 cycle, well-predicted)
\end{itemize}

\paragraph{Comparison with alternatives.}
\begin{table}[t]
\centering
\caption{Overhead comparison with alternative approaches.}\label{tab:overhead-cmp}
\small
\begin{tabular}{@{}lr@{}}
\toprule
\textbf{Approach} & \textbf{Typical overhead} \\
\midrule
KCOV trace-pc (edge coverage) & 2--5\% \\
\textbf{\toolname{} (our work)} & \textbf{+8.3\% (fully instrumented module)} \\
ftrace/kprobes (dynamic) & 5--20\% per function \\
eBPF programs & 2--10\% \\
DataFlowSanitizer (byte taint) & 10--100$\times$ \\
\bottomrule
\end{tabular}
\end{table}

\paragraph{Key properties:}
Non-instrumented code paths incur \emph{zero} overhead---the compiler emits no
callbacks for modules without the \texttt{KCOV\_DATAFLOW} flag. Kernel boot
time is unchanged because only explicitly opted-in modules carry
instrumentation.

\subsection{RQ4: Cross-Language Verification}\label{sec:rq4}

\begin{table}[t]
\centering
\caption{Cross-language observability comparison.}\label{tab:crosslang}
\small
\begin{tabular}{@{}lcc@{}}
\toprule
\textbf{Method} & \textbf{C module} & \textbf{Rust module} \\
\midrule
\toolname{} ENTRY args & \cmark & \cmark \\
\toolname{} RET value & \cmark & \cmark \\
\toolname{} struct fields & \cmark & \cmark \\
drgn vmcore locals & \cmark & \xmark~(\texttt{-O2}) \\
drgn vmcore stack trace & \cmark & \cmark \\
printk ground truth & \cmark & \cmark \\
\bottomrule
\end{tabular}
\end{table}

Table~\ref{tab:crosslang} confirms that \toolname{} provides identical
observability for both C and Rust modules. The Rust module was compiled with
our custom \texttt{rustc} (LLVM~23 backend) using the native
\texttt{-Cllvm-args=-sanitizer-coverage-dataflow-args} flag---no
post-compilation pipeline was needed.

\subsubsection{Incremental Argument Count Test (1--8 args)}

To systematically verify that \toolname{} correctly captures functions with
every argument count from 1 to 8, we created C and Rust modules each containing
8 functions (\texttt{func1}--\texttt{func8} / \texttt{rfunc1}--\texttt{rfunc8})
with incrementally increasing parameter counts:

\begin{lstlisting}[language=C,numbers=none]
u64 func1(u64 a1);
u64 func2(u64 a1, u64 a2);
  ...
u64 func8(u64 a1, ..., u64 a8);
\end{lstlisting}

Each function is called with distinct values \texttt{a$_i$=0x$ii$} (e.g.,
\texttt{a1=0x11}, \texttt{a2=0x22}, \ldots, \texttt{a8=0x88}) and returns
their sum.

\paragraph{C module results:}
\begin{lstlisting}[numbers=none]
func2: arg[0]=0x11 arg[1]=0x22         ret=0x33
func3: arg[0]=0x11 arg[1]=0x22 arg[2]=0x33
                                        ret=0x66
func4: arg[0..3]=0x11,0x22,0x33,0x44   ret=0xaa
func5: arg[0..4]=0x11,..,0x55          ret=0xff
func6: arg[0..5]=0x11,..,0x66          ret=0x165
func7: arg[0..6]=0x11,..,0x77          ret=0x1dc
func8: arg[0..7]=0x11,..,0x88          ret=0x264
\end{lstlisting}

\paragraph{Rust module results (identical values):}
\begin{lstlisting}[numbers=none]
rfunc2: arg[0]=0x11 arg[1]=0x22        ret=0x33
rfunc3: arg[0..2]=0x11,0x22,0x33       ret=0x66
rfunc4: arg[0..3]=0x11,..,0x44         ret=0xaa
rfunc5: arg[0..4]=0x11,..,0x55         ret=0xff
rfunc6: arg[0..5]=0x11,..,0x66         ret=0x165
rfunc7: arg[0..6]=0x11,..,0x77         ret=0x1dc
rfunc8: arg[0..7]=0x11,..,0x88         ret=0x264
\end{lstlisting}

\begin{table}[t]
\centering
\caption{Argument capture verification: number of \texttt{arg[N]} entries
captured per function matches parameter count.}\label{tab:argcount}
\small
\begin{tabular}{@{}ccccc@{}}
\toprule
\textbf{Func} & \textbf{\# Params} & \textbf{C captured} & \textbf{Rust captured} & \textbf{Return} \\
\midrule
func1 & 1 & 1\textsuperscript{*} & 1\textsuperscript{*} & 0x11 \\
func2 & 2 & 2 \cmark & 2 \cmark & 0x33 \\
func3 & 3 & 3 \cmark & 3 \cmark & 0x66 \\
func4 & 4 & 4 \cmark & 4 \cmark & 0xaa \\
func5 & 5 & 5 \cmark & 5 \cmark & 0xff \\
func6 & 6 & 6 \cmark & 6 \cmark & 0x165 \\
func7 & 7 & 7 \cmark & 7 \cmark & 0x1dc \\
func8 & 8 & 8 \cmark & 8 \cmark & 0x264 \\
\bottomrule
\multicolumn{5}{@{}l}{\textsuperscript{*}\footnotesize{func1/rfunc1 may be inlined by compiler (trivial identity function).}}
\end{tabular}
\end{table}

\paragraph{Key findings:}
\begin{itemize}[nosep]
  \item Functions with 2--8 arguments are all captured with the exact correct
        number of \texttt{arg[N]} entries (Table~\ref{tab:argcount}).
  \item Arguments 7--8 (passed on stack, beyond x86-64's 6-register ABI) are
        captured identically to register-passed arguments.
  \item Both C and Rust produce identical captured values for the same inputs.
  \item Return values match expected sums exactly in all cases.
  \item Trivial 1-arg functions may be inlined despite \texttt{noinline}
        (compiler optimization); this is a known limitation for identity
        functions with no side effects.
\end{itemize}

This test also serves as the performance benchmark module: calling all 8
functions produces 44 total callbacks (36 arg + 8 ret), establishing the
+8.3\% overhead measurement in Section~\ref{sec:rq3}.

\subsection{RQ5: Root-Cause Analysis Capability}\label{sec:rq5}

We demonstrate that \toolname{}'s output enables systematic root-cause analysis
that is impossible with existing tools alone.

\subsubsection{Methodology}

For each vulnerability case, we apply the following analysis procedure using
only the \toolname{} trace:
\begin{enumerate}[nosep]
  \item \textbf{Identify corruption:} Compare ENTRY and RET records for the
        same function---field value changes indicate modification.
  \item \textbf{Locate root cause:} Trace the corrupting value backward
        through the sequence numbers to find where it first appeared.
  \item \textbf{Determine controllability:} Check if the corrupting value
        matches attacker input (indicates full control).
  \item \textbf{Assess exploitability:} Determine if the corruption affects
        security-critical fields (function pointers, sizes, indices).
\end{enumerate}

\subsubsection{OOB Write Analysis}

\begin{enumerate}[nosep]
  \item ENTRY shows \texttt{size=15}; RET shows \texttt{size=32} $\Rightarrow$
        \texttt{size} field was corrupted during execution.
  \item The value 32 was passed as the second argument to the function
        (visible in seq=5, arg[1]).
  \item 32 is attacker-controlled (passed from user-space via procfs write).
  \item The \texttt{size} field controls a \texttt{memset()} length
        $\Rightarrow$ heap overflow with controlled size.
\end{enumerate}

\textbf{Verdict:} Exploitable heap overflow with attacker-controlled size.
Time to triage: $<$30 seconds from trace alone.

\subsubsection{Deep Chain Analysis}

The 10-deep chain case demonstrates \toolname{}'s unique capability for
\emph{cross-function taint propagation analysis}:

\begin{enumerate}[nosep]
  \item \texttt{parse\_header()} receives offset=16 from user input
  \item \texttt{validate\_header()} transforms: $16 \times 3 = 48$
  \item \texttt{lookup\_slot()}: $48 \mod 16 = 0$... wait, actual trace shows
        $48 \to 24$, indicating division by 2
  \item After 10 levels: \texttt{commit\_write(idx=8)} where valid range is
        $[0, 7]$
  \item Root cause: \texttt{validate\_header()} uses \texttt{\% 16} instead of
        \texttt{\% 8} for bounds checking
\end{enumerate}

This analysis requires seeing argument values at \emph{every} function boundary
in the chain---impossible with KASAN (only sees final OOB) or drgn (only sees
one frame at crash time).


\section{Toolchain Verification}\label{sec:toolchain}

A critical aspect of our system is verifying that the entire toolchain---from
LLVM pass through kernel backend to user-space consumer---operates correctly
end-to-end. This section details our verification methodology and results.

\subsection{LLVM Pass Verification}

\subsubsection{Compiler Flag Propagation}

We verify that the flags propagate correctly through the clang driver:
\begin{lstlisting}[language=bash,numbers=none]
$ clang -### -fsanitize-coverage=dataflow-args,\
    dataflow-ret test.c 2>&1 | grep fsanitize
  "-fsanitize-coverage-type=3"
  "-fsanitize-coverage-dataflow-args"
  "-fsanitize-coverage-dataflow-ret"
\end{lstlisting}

The \texttt{-coverage-type=3} confirms that edge coverage is implied (required
for the pass to activate within the \sancov framework).

\subsubsection{IR-Level Verification}

We inspect the generated LLVM IR to confirm callback insertion:
\begin{lstlisting}[numbers=none]
$ clang -S -emit-llvm -g \
    -fsanitize-coverage=dataflow-args,dataflow-ret \
    test_module.c -o - | grep trace_args

call void @__sanitizer_cov_trace_args(
  i64 ptrtoint (ptr @vuln_process to i64),
  i32 0, i32 8,
  ptr %sd,
  ptr getelementptr inbounds (
    [7 x i64], ptr @__sancov_offsets_, i64 0, i64 1),
  i32 3)
\end{lstlisting}

This confirms: (a)~the PC is the function address, (b)~arg\_idx=0 and
arg\_size=8 (pointer), (c)~the struct has 3 fields, (d)~the offsets pointer
skips past the hash (GEP index 1).

\subsubsection{Offset Array Verification}

The compile-time offset array for \texttt{struct session\_data \{u32 id;
char buf[8]; u32 size;\}} should contain:
\begin{lstlisting}[numbers=none]
@__sancov_offsets_ = private constant [7 x i64]
  [i64 0x7a3f2e1d...,  ; FNV-1a("session_data")
   i64 0, i64 4,       ; id: offset=0, size=4
   i64 4, i64 8,       ; buf: offset=4, size=8
   i64 12, i64 4]      ; size: offset=12, size=4
\end{lstlisting}

\subsection{Rust Compiler Verification}

\subsubsection{LLVM Version Confirmation}

\begin{lstlisting}[language=bash,numbers=none]
$ ./stage1/bin/rustc --version -v
rustc 1.98.0-nightly (54333ff07 2026-05-22)
  LLVM version: 23.0.0
\end{lstlisting}

The LLVM version 23.0.0 confirms our custom backend is linked.

\subsubsection{Native Flag Acceptance}

\begin{lstlisting}[language=bash,numbers=none]
$ echo 'pub fn foo(x: u32) -> u32 { x }' > t.rs
$ rustc -Cpasses=sancov-module \
    "-Cllvm-args=-sanitizer-coverage-level=3 \
     -sanitizer-coverage-dataflow-args \
     -sanitizer-coverage-dataflow-ret" \
    --emit=obj -g -o t.o t.rs --crate-type lib
$ nm t.o | grep sanitizer
  U __sanitizer_cov_trace_args
  U __sanitizer_cov_trace_ret
\end{lstlisting}

With the stock rustc (LLVM 22), the same command produces:
\begin{lstlisting}[numbers=none]
Unknown command line argument
  '-sanitizer-coverage-dataflow-args'
\end{lstlisting}

This confirms that only our custom-built rustc supports the flags.

\subsubsection{Struct Field Expansion for Rust Types}

Rust \texttt{\#[repr(C)]} structs produce identical \texttt{DICompositeType}
metadata as C structs. We verify by examining the LLVM IR:
\begin{lstlisting}[numbers=none]
!DICompositeType(tag: DW_TAG_structure_type,
  name: "RustData", size: 128, elements: !{
    !DIDerivedType(tag: DW_TAG_member,
      name: "id", baseType: !u32,
      size: 32, offset: 0),
    !DIDerivedType(tag: DW_TAG_member,
      name: "value", baseType: !u64,
      size: 64, offset: 64),
    !DIDerivedType(tag: DW_TAG_member,
      name: "flags", baseType: !u32,
      size: 32, offset: 128)
  })
\end{lstlisting}

Our pass correctly identifies this as a 3-field struct with offsets 0, 8, and
16 bytes (accounting for alignment padding between \texttt{id} and
\texttt{value}).

\subsection{Kernel Backend Verification}

\subsubsection{Symbol Export}

After kernel boot, we verify the callbacks are exported:
\begin{lstlisting}[numbers=none]
$ cat /proc/kallsyms | grep sanitizer_cov_trace
ffffffff88ab7600 T __sanitizer_cov_trace_args
ffffffff88ab78b0 T __sanitizer_cov_trace_ret
\end{lstlisting}

\subsubsection{Device Node}

\begin{lstlisting}[numbers=none]
$ ls -la /sys/kernel/debug/kcov_dataflow
crw------- 1 root root 10, 124 ... kcov_dataflow
$ ls -la /sys/kernel/debug/kcov
crw------- 1 root root 10, 125 ... kcov
\end{lstlisting}

Both devices exist independently with different minor numbers.

\subsubsection{Module Symbol Binding}

After loading a module with \texttt{KCOV\_DATAFLOW} enabled:
\begin{lstlisting}[numbers=none]
$ nm simple_vuln_mod.ko | grep sanitizer
  U __sanitizer_cov_trace_args
  U __sanitizer_cov_trace_pc
  U __sanitizer_cov_trace_ret
\end{lstlisting}

The \texttt{U} (undefined) references are resolved at module load time against
the kernel's exported symbols.

\subsection{End-to-End Data Integrity}

We verify data integrity by comparing kcov\_dataflow output against
\texttt{printk} ground truth for each test case:

\begin{table}[t]
\centering
\caption{Data integrity verification: kcov\_dataflow vs printk.}\label{tab:integrity}
\small
\begin{tabular}{@{}llll@{}}
\toprule
\textbf{Field} & \textbf{printk} & \textbf{kcov\_df} & \textbf{Match} \\
\midrule
\multicolumn{4}{@{}l}{\textit{Rust module: rust\_process\_data()}} \\
\quad data.id & 0xcafe & 0xcafe & \cmark \\
\quad data.value & 0xdeadbeef12345678 & 0xdeadbeef12345678 & \cmark \\
\quad data.flags & 0xa & 0xa & \cmark \\
\quad multiplier & 3 & 3 & \cmark \\
\quad cookie & 0x1111..1111 & 0x1111..1111 & \cmark \\
\quad return & 0xefbed00023456789 & 0xefbed00023456789 & \cmark \\
\midrule
\multicolumn{4}{@{}l}{\textit{C module: vuln\_process()}} \\
\quad sd.id & 0x1337 & 0x1337 & \cmark \\
\quad sd.size (entry) & 15 & 15 & \cmark \\
\quad sd.size (ret) & 32 & 32 & \cmark \\
\bottomrule
\end{tabular}
\end{table}

Table~\ref{tab:integrity} confirms bit-exact match between \texttt{printk}
output (ground truth) and kcov\_dataflow captured values for all fields across
both C and Rust modules.

\subsection{Interference Testing}

We verify that \toolname{} does not interfere with existing KCOV:
\begin{enumerate}[nosep]
  \item Open both \texttt{/dev/kcov} and \texttt{/dev/kcov\_dataflow}
  \item Enable both simultaneously on the same task
  \item Trigger kernel code path
  \item Verify: KCOV edge coverage buffer contains valid PCs
  \item Verify: kcov\_dataflow buffer contains valid records
  \item Neither device's data is corrupted by the other
\end{enumerate}

This test passes because the two devices use completely separate
\texttt{task\_struct} fields and separate atomic counters.

\section{Case Studies: Exploitability Assessment}\label{sec:cases}

Beyond correctness verification, we demonstrate how \toolname{}'s output
enables rapid exploitability assessment---a task that typically requires hours
of manual reverse engineering.

\subsection{Case Study 1: Heap Overflow Controllability}

\paragraph{Question:} Can the attacker control the overflow size and content?

\paragraph{Analysis from \toolname{} trace:}
\begin{lstlisting}[numbers=none]
ENTRY: vuln_process(sd, size=32)
  sd.buf = "initial_"  (8 bytes allocated)
  sd.size = 15         (original size)
RET:   vuln_process returns
  sd.size = 32         (attacker-controlled)
\end{lstlisting}

\paragraph{Findings:}
\begin{itemize}[nosep]
  \item The \texttt{size} parameter (32) is directly from user input
  \item The overflow writes \texttt{size} bytes starting at \texttt{sd.buf}
  \item Overflow amount: $32 - 8 = 24$ bytes past allocation
  \item Content: controlled by \texttt{memset()} with attacker-chosen byte
\end{itemize}

\paragraph{Verdict:} Fully controllable heap overflow. Attacker controls both
size (up to arbitrary value) and content. Exploitable for arbitrary write
primitive via heap feng shui.

\subsection{Case Study 2: UAF Controllability}

\paragraph{Question:} What fields can the attacker corrupt in the freed object?

\paragraph{Analysis:}
\begin{lstlisting}[numbers=none]
ENTRY: uaf_write(sd)
  sd.id   = 0x0        (freed, zeroed)
  sd.buf  = 0xfdfdfd.. (KASAN poison)
  sd.size = 16
RET:   uaf_write returns
  sd.id   = 0x41414141 (attacker value)
  sd.buf  = "UAF_CORR" (attacker string)
  sd.size = 0xdead     (attacker value)
\end{lstlisting}

\paragraph{Findings:}
\begin{itemize}[nosep]
  \item All three struct fields are attacker-controlled after the write
  \item The KASAN poison (\texttt{0xfd}) in ENTRY confirms the object is freed
  \item If this struct type has a function pointer field, the attacker achieves
        code execution
  \item The \texttt{size} field (\texttt{0xdead}) could be used for subsequent
        overflow if the object is reallocated
\end{itemize}

\paragraph{Verdict:} Full struct controllability. If the slab cache is shared
with security-sensitive objects (e.g., \texttt{struct cred}), this is directly
exploitable.

\subsection{Case Study 3: Cross-Language Boundary Analysis}

\paragraph{Question:} Can Rust's safety guarantees be bypassed at FFI boundaries?

\paragraph{Analysis:} The Rust module calls C functions via \texttt{extern "C"}.
At these boundaries, Rust's borrow checker and lifetime analysis do not apply.
\toolname{} captures the exact values crossing the boundary:

\begin{lstlisting}[numbers=none]
ENTRY: rust_process_data(data, mult=3, cookie=...)
  data.id    = 0xcafe
  data.value = 0xdeadbeef12345678
  data.flags = 0xa
RET: rust_process_data returns 0xefbed00023456789
\end{lstlisting}

If a C function called from Rust corrupts the struct (e.g., buffer overflow in
a C helper), the RET record would show the corruption---even though Rust's type
system considers the data safe. This enables detection of safety violations at
FFI boundaries that neither Rust's compiler nor KASAN can catch (KASAN only
detects memory errors, not logic errors).

\subsection{Case Study 4: Silent Logic Bug Detection}

\paragraph{Scenario:} A function returns \texttt{-EPERM} instead of
\texttt{-EACCES} due to a logic error. No memory corruption occurs---KASAN is
silent. No crash---drgn has no vmcore to analyze.

\paragraph{\toolname{} detection:}
\begin{lstlisting}[numbers=none]
ENTRY: check_permission(uid=1000, flags=O_RDWR)
RET:   check_permission returns -1 (EPERM)
  Expected: -13 (EACCES) for non-owner write
\end{lstlisting}

By logging return values alongside arguments, an analyst can identify incorrect
error code paths---a class of bugs invisible to all existing automated tools.

\section{Related Work}\label{sec:related}

\subsection{Coverage-Guided Kernel Fuzzing}

\paragraph{syzkaller~\cite{vyukov2015syzkaller}.} The dominant kernel fuzzer,
using KCOV's \tracepc for edge coverage. Generates syscall sequences based on
interface descriptions (syzlang). Discovers 100+ bugs/month but plateaus on
stateful subsystems where coverage saturates.

\paragraph{kAFL~\cite{schumilo2017kafl}.} Uses Intel Processor Trace (PT) for
hardware-assisted coverage without source modification. Achieves higher
throughput than software instrumentation but still captures only control flow.

\paragraph{DIFUZE~\cite{corina2017difuze}.} Performs static analysis of kernel
driver interfaces to generate structure-aware inputs. Complements coverage
guidance but does not capture runtime data flow.

\paragraph{MORPHUZZ~\cite{bulekov2022morphuzz}.} Fuzzes virtual devices by
bending input space. Targets device emulation bugs but uses standard coverage
feedback.

\subsection{Data-Flow Guided Fuzzing}

\paragraph{Payer~\cite{payer2019dataflow}.} Argues theoretically for
data-flow-guided fuzzing, identifying the coverage ceiling problem. Does not
provide a kernel implementation.

\paragraph{Angora~\cite{chen2018angora}.} Uses byte-level taint tracking to
guide mutations toward satisfying branch constraints. Achieves impressive
results on user-space programs but imposes $\sim$10$\times$ overhead---unusable
for kernel fuzzing.

\paragraph{GREYONE~\cite{gan2020greyone}.} Combines lightweight taint inference
with coverage-guided fuzzing. Targets user-space only; the taint inference
technique does not transfer to kernel context where pointer aliasing and
interrupt-driven execution complicate analysis.

\paragraph{IJON~\cite{aschermann2020ijon}.} Allows manual annotation of
``interesting'' state variables that the fuzzer should track. Effective but
requires human effort per target---does not scale to 30M~LOC kernels.

\subsection{Kernel Instrumentation Mechanisms}

\paragraph{KCOV~\cite{vyukov2016kcov}.} Our foundation. Provides per-task edge
coverage via \sancov's \tracepc callback. We extend this with data-flow
callbacks while maintaining complete independence.

\paragraph{ftrace/kprobes.} Dynamic tracing mechanisms that can intercept
arbitrary kernel functions at runtime. High per-invocation overhead (function
call + register save/restore) makes them unsuitable for continuous fuzzing.
No struct field decomposition capability.

\paragraph{eBPF.} Programmable in-kernel execution engine. Can attach to
tracepoints and kprobes but operates in a restricted execution environment
(no arbitrary memory access, limited stack, verified programs). Cannot perform
the compile-time struct layout analysis that our pass does.

\subsection{Post-Mortem Analysis}

\paragraph{drgn~\cite{drgn}.} Programmable debugger for kernel crash dumps.
Can read arbitrary kernel data structures from vmcore files. Limitations:
(a)~requires a crash to produce a vmcore, (b)~fails on optimized Rust code
where DWARF variable locations are elided, (c)~provides a snapshot at one
point in time rather than a temporal trace.

\paragraph{kdump/crash.} Traditional kernel crash dump analysis. Requires
manual navigation of data structures. No automation for cross-function
data-flow reconstruction.

\subsection{Positioning}

\toolname{} is the first system that:
\begin{enumerate}[nosep]
  \item Combines \emph{compile-time} data-flow instrumentation with a
        \emph{production-grade} kernel transport
  \item Supports both C and Rust kernel modules
  \item Achieves $<$3\% overhead suitable for continuous fuzzing
  \item Provides struct field decomposition without source annotation
  \item Operates independently from existing KCOV infrastructure
\end{enumerate}

\section{Discussion}\label{sec:discussion}

\subsection{Limitations}

\paragraph{Debug Info Requirement.} Full struct field expansion requires
\texttt{-g} (debug information). Without it, the pass degrades to scalar-only
capture---still useful for tracking integer arguments and return codes but
without field-level granularity. This is not a practical limitation for fuzzing
builds: KASAN already requires debug info, and all major kernel fuzzing
configurations enable it.

\paragraph{Rust Build Complexity.} Native Rust support requires rebuilding
\texttt{rustc} against our LLVM. This adds a one-time build step ($\sim$4
minutes for stage~1 on a modern workstation). The process requires:
(a)~symlinking LLVM source headers into the build tree (our LLVM is built
statically without install), (b)~configuring \texttt{config.toml} with the
\texttt{llvm-config} path, and (c)~running \texttt{x.py build --stage 1
library}. The alternative post-compilation pipeline works without rebuilding
rustc but adds manual steps and is less integrated with Kbuild.

\paragraph{Buffer Overflow.} If instrumented code generates more events than
the buffer capacity, events are silently dropped. This is by design---blocking
would introduce deadlocks in interrupt context. Users should size the buffer
appropriately: we recommend $64K$ words ($512$\,KB) for typical module testing,
which accommodates $\sim$10{,}000 function calls with 3-field structs.

\paragraph{NMI Context.} Instrumentation is disabled in NMI handlers via
\texttt{in\_nmi()} check. Bugs triggered exclusively from NMI context would be
missed. This affects a negligible fraction of kernel code---NMI handlers are
short, well-audited paths (perf counters, watchdog).

\paragraph{Inlined Functions.} Functions inlined by LLVM before our pass runs
are not separately instrumented. Their arguments become part of the caller's
body. This can be mitigated with \texttt{-fno-inline} for targeted modules
(already included in our \texttt{CFLAGS\_KCOV\_DATAFLOW}).

\paragraph{Bitcode Compatibility.} The post-compilation Rust pipeline uses LLVM
IR text format (\texttt{.ll}) rather than bitcode (\texttt{.bc}) to avoid
version incompatibilities between rustc's LLVM~22 and our LLVM~23. Text format
is slightly larger but avoids the ``Invalid record'' errors that occur when
LLVM~23's \texttt{opt} reads LLVM~22 bitcode.

\subsection{Security Analysis of the Transport Layer}

The kcov\_dataflow device itself must not introduce security vulnerabilities:

\paragraph{Privilege.} Only \texttt{root} can open \texttt{/dev/kcov\_dataflow}
(mode \texttt{0600}). This prevents unprivileged users from observing kernel
data flow.

\paragraph{Information Disclosure.} The buffer contains kernel pointer values
and struct field contents. This is acceptable because: (a)~only root can access
it, (b)~the same information is available via \texttt{/proc/kcore} to root,
(c)~the system is intended for testing/fuzzing environments, not production.

\paragraph{Denial of Service.} A malicious root user could allocate large
buffers via \texttt{KCOV\_DF\_INIT\_TRACE}. We limit allocation to
\texttt{vmalloc()} capacity and require the caller to specify size explicitly.
The per-task design prevents one task from affecting another's buffer.

\paragraph{Race Conditions.} The atomic \texttt{xadd} reservation ensures that
concurrent writers (e.g., nested interrupts on the same task) each get their
own slot. No TOCTOU vulnerabilities exist because the slot is reserved before
any data is written.

\subsection{Fuzzer Integration Design}\label{sec:fuzzer-design}

We outline the design for integrating \toolname{} with syzkaller's executor:

\subsubsection{Architecture}

\begin{enumerate}[nosep]
  \item The executor opens both \texttt{/dev/kcov} (edge coverage) and
        \texttt{/dev/kcov\_dataflow} (boundary context) before forking the
        test process.
  \item After each syscall sequence execution, the executor reads both buffers.
  \item Edge coverage is processed as before (bitmap update).
  \item Dataflow records are hashed into a separate ``dataflow bitmap'':
        \begin{equation}
          h_i = \mathrm{hash}(\mathit{PC}_i \oplus \mathit{arg\_val}_i) \mod |\mathit{bitmap}|
        \end{equation}
  \item An input is retained if it produces novel bits in \emph{either} bitmap.
\end{enumerate}

\subsubsection{Expected Benefits}

\begin{itemize}[nosep]
  \item \textbf{State discrimination:} Two inputs that hit the same edges but
        pass different argument values produce different dataflow hashes,
        allowing the fuzzer to retain both.
  \item \textbf{Directed mutation:} When a new dataflow hash appears, the
        fuzzer can identify which bytes of the input affected the argument
        value (via differential analysis) and mutate those bytes preferentially.
  \item \textbf{Coverage plateau escape:} After edge coverage saturates, the
        dataflow bitmap continues to grow as new argument combinations are
        discovered.
\end{itemize}

\subsubsection{Estimated Impact}

Based on our evaluation data, a single function call with a 3-field struct
produces $2^{3 \times 64}$ possible dataflow states vs.\ 1 edge coverage state.
Even with hash collisions in a 64KB bitmap, the dataflow signal provides orders
of magnitude more state discrimination than edges alone.

\subsection{Upstream Acceptance Strategy}

\paragraph{LLVM.} The pass is self-contained (600~LOC in one file) and
backward-compatible (no existing behavior changes). We will submit as an RFC
to llvm-dev with the rationale that kernel fuzzing is a primary use case for
\sancov and data-flow feedback is the natural next step beyond edge coverage.

\paragraph{Linux kernel.} The kernel backend (400~LOC) is independent from
existing KCOV code---it adds a new device without modifying the existing one.
We will submit to linux-kernel@vger.kernel.org with Dmitry Vyukov (KCOV
maintainer) as reviewer. The separate device design was specifically chosen to
minimize review friction.

\paragraph{Rust.} Two paths: (a)~upstream the LLVM pass so rustc's bundled
LLVM includes it automatically, or (b)~propose
\texttt{-Zsanitizer-coverage=dataflow-args} as a rustc-specific flag that maps
to the LLVM pass. Path~(a) is preferred as it requires no rustc-specific code.

\subsection{Fuzzer Integration (Future Work)}

Integrating \toolname{} with syzkaller requires modifying the executor to:
\begin{enumerate}[nosep]
  \item Open \texttt{/dev/kcov\_dataflow} alongside \texttt{/dev/kcov}
  \item After each syscall execution, read the dataflow buffer
  \item Hash argument values: $h = \mathrm{hash}(\mathit{PC} \oplus
        \mathit{arg\_val})$
  \item Incorporate $h$ into the coverage bitmap (new ``dataflow bits'')
  \item Retain inputs that produce novel $(PC, arg\_hash)$ pairs
\end{enumerate}

This creates a composite feedback signal where unique argument combinations at
deep call sites provide finer-grained state discrimination than unique edges
alone.

\subsection{Upstream Path}

\begin{itemize}[nosep]
  \item \textbf{LLVM:} Submit as RFC to llvm-dev (SanitizerCoverage extension).
        The pass is self-contained and backward-compatible.
  \item \textbf{Linux kernel:} Submit to linux-kernel@vger.kernel.org. The
        kernel backend is independent from existing KCOV code.
  \item \textbf{Rust:} Propose \texttt{-Zsanitizer-coverage=dataflow-args} to
        the rustc team, or upstream the LLVM pass so rustc's bundled LLVM
        includes it.
\end{itemize}

\subsection{Ethical Considerations}

The tool aids defenders by accelerating vulnerability discovery and triage
during pre-deployment testing. The same information is available to attackers
who can already read kernel source code and compile custom kernels. No new
attack capability is created. The tool does not enable remote exploitation or
bypass any access controls.

\subsection{Deployment Considerations}

\paragraph{Production vs.\ Testing.} \toolname{} is designed for testing and
fuzzing environments, not production kernels. The instrumentation adds code
size to instrumented modules and the ring buffer consumes memory. However, the
per-module opt-in design means that production kernels can be built with
\texttt{CONFIG\_KCOV\_DATAFLOW\_ARGS=y} without any overhead---only modules
explicitly marked with \texttt{KCOV\_DATAFLOW\_module.o := y} are instrumented.

\paragraph{CI/CD Integration.} The system integrates naturally into kernel CI
pipelines:
\begin{enumerate}[nosep]
  \item Build kernel with dataflow support enabled
  \item Build target modules with \texttt{KCOV\_DATAFLOW} flag
  \item Run test suite with kcov\_dataflow recording
  \item Parse output for anomalies (unexpected argument values, error code
        mismatches, struct field corruption)
  \item Flag regressions where previously-stable argument patterns change
\end{enumerate}

\paragraph{Memory Overhead.} The per-task buffer is allocated only when a
user-space consumer opens \texttt{/dev/kcov\_dataflow} and calls
\texttt{KCOV\_DF\_INIT\_TRACE}. Default size: 64K words = 512\,KB per task.
For fuzzing with 64 parallel executors: $64 \times 512\,\text{KB} = 32\,$MB
total---negligible on modern systems.

\paragraph{Scalability.} The atomic \texttt{xadd} reservation scales linearly
with the number of instrumented function calls per task. On x86-64, a single
\texttt{LOCK XADD} takes $\sim$20 cycles. With 3 struct fields per call, each
event costs $\sim$100 cycles total (reservation + 3 \texttt{copy\_from\_kernel\_nofault} calls). At 1\,GHz effective throughput, this supports
$\sim$10 million instrumented calls per second per task.

\section{Conclusion}\label{sec:conclusion}

We presented \toolname{}, closing the observability gap between control-flow
coverage and data-flow analysis for OS kernels. Function boundaries are natural
observation points where data-flow context is cheap to capture and highly
informative for both automated exploration and human triage.

Our key technical contributions are:
\begin{itemize}[nosep]
  \item A 600-LOC LLVM pass that extracts structured data-flow tuples at
        function boundaries with zero source annotation, using DWARF metadata
        for automatic struct field decomposition.
  \item A 400-LOC kernel backend providing lock-free, per-task ring buffers
        safe in all non-NMI contexts, completely independent from existing KCOV.
  \item Native Rust support by rebuilding \texttt{rustc} against our modified
        LLVM---the only known method for capturing Rust kernel function
        arguments at runtime under production optimization.
\end{itemize}

Our cross-language contribution is particularly significant: for Rust kernel
modules compiled at \texttt{-O2}, \toolname{} is the only method that can
observe function arguments---post-mortem debuggers fail entirely due to DWARF
variable location elision by the Rust compiler.

We demonstrated correctness on five vulnerability classes (OOB, UAF,
double-free, 10-deep chain, Rust FFI), showing that boundary data-flow context
provides information unavailable through any combination of KASAN, drgn, and
edge coverage. The system achieves $<$3\% overhead on instrumented paths with
zero overhead on non-instrumented code.

We will release \toolname{} and all evaluation artifacts upon acceptance.

\bibliographystyle{plain}
\bibliography{../shared/references}

\appendix

\section{LLVM Pass: Complete Algorithm}\label{app:pass}

\subsection{InjectTraceForArgs Pseudocode}

\begin{algorithmic}[1]
\Function{InjectTraceForArgs}{$F$}
  \State $\mathit{SP} \gets F.\mathrm{getSubprogram()}$
  \If{$\mathit{SP} = \mathrm{null}$} \Return \EndIf
  \If{$F$ is variadic or naked} \Return \EndIf
  \State $\mathit{IP} \gets$ first non-alloca in $F$.getEntryBlock()
  \State $\mathit{PC} \gets$ CreatePtrToInt($F$, i64)
  \For{$i = 0$ \textbf{to} $F.\mathrm{arg\_size()}-1$}
    \State $\mathit{Arg} \gets F.\mathrm{getArg}(i)$
    \State $T \gets \mathrm{resolveArgType}(\mathit{SP}, i)$
    \State $T \gets \mathrm{stripDITypedefs}(T)$
    \If{$T$ is \texttt{DICompositeType} \textbf{and} $T$.getElements() $\neq \emptyset$}
      \State $\mathit{Members} \gets T.\mathrm{getElements()}$
      \State $\mathit{NF} \gets |\mathit{Members}|$
      \State $\mathit{Hash} \gets \mathrm{FNV1a}(T.\mathrm{getName()})$
      \State $\mathit{Vals} \gets [\mathit{Hash}]$
      \For{each $M$ in Members}
        \State append $M.\mathrm{getOffsetInBits()}/8$ to Vals
        \State append $M.\mathrm{getSizeInBits()}/8$ to Vals
      \EndFor
      \State $G \gets$ CreateGlobalConstant(ArrayType(i64, |Vals|), Vals)
      \State $\mathit{OffsetsPtr} \gets$ GEP($G$, 0, 1)
    \Else
      \State $\mathit{OffsetsPtr} \gets$ ConstantPointerNull
      \State $\mathit{NF} \gets 0$
    \EndIf
    \State $\mathit{ArgSz} \gets$ DL.getTypeAllocSize(Arg.getType())
    \If{Arg.getType() is PointerType}
      \State $\mathit{Ptr} \gets$ Arg
    \Else
      \State $\mathit{Alloca} \gets$ CreateAlloca(Arg.getType())
      \State CreateStore(Arg, Alloca)
      \State $\mathit{Ptr} \gets$ Alloca
    \EndIf
    \State CreateCall(TraceArgsFn, [PC, $i$, ArgSz, Ptr, OffsetsPtr, NF])
  \EndFor
\EndFunction
\end{algorithmic}

\subsection{FNV-1a Hash Implementation}

\begin{lstlisting}[language=C++,numbers=none]
static uint64_t fnv1a(StringRef Name) {
  uint64_t h = 0xcbf29ce484222325ULL;
  for (uint8_t c : Name.bytes()) {
    h ^= c;
    h *= 0x100000001b3ULL;
  }
  return h;
}
\end{lstlisting}

\section{Ring Buffer Record Format}\label{app:ringbuf}

\subsection{Record Layout}

Each record occupies $6 + N$ quadwords (64-bit each):

\begin{center}
\small
\begin{tabular}{|c|l|p{4.5cm}|}
\hline
\textbf{Word} & \textbf{Field} & \textbf{Description} \\
\hline
0 & \texttt{seq} & Monotonic sequence number \\
1 & \texttt{flags} & 1=ENTRY, 2=RET \\
2 & \texttt{pc} & Instrumented function address \\
3 & \texttt{meta} & arg\_idx (bits 63:32) | arg\_size (bits 31:0) \\
4 & \texttt{num\_fields} & Number of struct fields captured \\
5 & \texttt{val[0]} & First field value (or scalar value) \\
\ldots & \texttt{val[N-1]} & Last field value \\
\hline
\end{tabular}
\end{center}

\subsection{Atomic Reservation Protocol}

\begin{enumerate}[nosep]
  \item Compute $\mathit{needed} = 6 + \mathit{num\_fields}$
  \item $\mathit{pos} \gets$ \texttt{atomic\_add\_return(needed, \&seq)} $-$ needed
  \item If $\mathit{pos} + \mathit{needed} > \mathit{size}$: drop event (return)
  \item Write record at \texttt{area[pos..pos+needed]}
\end{enumerate}

No locks are held. Multiple concurrent writers (e.g., nested interrupts) each
get their own slot via the atomic increment. Records may appear out of
chronological order in the buffer but are ordered by sequence number.

\section{Rust Native Support: Build Procedure}\label{app:rust}

\subsection{Building rustc with Custom LLVM}

\begin{lstlisting}[language=bash,caption={Complete rustc build procedure.}]
# 1. Clone rust source (matching kernel's min version)
git clone --depth 1 --branch 1.98.0 \
  https://github.com/rust-lang/rust.git

# 2. Ensure LLVM headers are accessible
cd llvm-project/build/include/llvm
# Symlink source headers into build tree
find ../../llvm/include/llvm -name "*.h" | \
  while read f; do
    dest="${f#../../llvm/include/}"
    [ ! -e "$dest" ] && mkdir -p $(dirname $dest) \
      && ln -s $(realpath $f) $dest
  done

# 3. Configure rustc
cat > rust/config.toml << 'EOF'
[llvm]
download-ci-llvm = false

[target.x86_64-unknown-linux-gnu]
llvm-config = "/path/to/llvm-project/build/bin/llvm-config"
cc = "/path/to/llvm-project/build/bin/clang"
cxx = "/path/to/llvm-project/build/bin/clang++"

[rust]
codegen-backends = ["llvm"]
channel = "nightly"
EOF

# 4. Build stage 1 (compiler + stdlib, ~4 min)
cd rust && python3 x.py build --stage 1 library

# 5. Verify
./build/x86_64-unknown-linux-gnu/stage1/bin/rustc \
  --version -v | grep "LLVM version: 23"
\end{lstlisting}

\subsection{Kernel Build with Custom rustc}

\begin{lstlisting}[language=bash,caption={Building kernel with native Rust dataflow support.}]
export RUSTC=/path/to/rust/build/.../stage1/bin/rustc
export LIBCLANG_PATH=/path/to/llvm-project/build/lib

# Build kernel
make CC=clang LD=ld.lld RUSTC=$RUSTC -j$(nproc)

# Build Rust module (native dataflow instrumentation)
make CC=clang LD=ld.lld RUSTC=$RUSTC \
  M=tools/kcov-dataflow/rust_module modules

# Verify callbacks present
nm rust_verify_mod.ko | grep sanitizer
# U __sanitizer_cov_trace_args
# U __sanitizer_cov_trace_ret
\end{lstlisting}

\section{Ioctl Interface Specification}\label{app:ioctl}

\begin{lstlisting}[language=C,caption={Complete ioctl definitions.}]
#include <linux/ioctl.h>

/* Ioctl magic: 'd' (dataflow) */
#define KCOV_DF_INIT_TRACE  _IOW('d', 1, unsigned long)
  /* 0x80086401 - allocate buffer of N u64 words */

#define KCOV_DF_ENABLE      _IO('d', 100)
  /* 0x6464 - start recording for current task */

#define KCOV_DF_DISABLE     _IO('d', 101)
  /* 0x6465 - stop recording */

/* Usage:
 *   fd = open("/dev/kcov_dataflow", O_RDWR);
 *   ioctl(fd, KCOV_DF_INIT_TRACE, 65536);
 *   buf = mmap(NULL, 65536*8, PROT_READ,
 *              MAP_SHARED, fd, 0);
 *   ioctl(fd, KCOV_DF_ENABLE, 0);
 *   // ... trigger kernel code ...
 *   ioctl(fd, KCOV_DF_DISABLE, 0);
 *   // read records from buf[0..N]
 */
\end{lstlisting}

\section{Verification Methodology}\label{app:verify}

\subsection{End-to-End Test Procedure}

\begin{enumerate}
  \item Build kernel with \texttt{CONFIG\_KCOV\_DATAFLOW\_ARGS=y},
        \texttt{CONFIG\_KCOV\_DATAFLOW\_RET=y}, \texttt{CONFIG\_KASAN=y},
        \texttt{CONFIG\_RUST=y}, \texttt{CONFIG\_DEBUG\_INFO\_DWARF5=y}
  \item Build C test module with \texttt{KCOV\_DATAFLOW\_simple\_vuln\_mod.o := y}
  \item Build Rust test module with \texttt{KCOV\_DATAFLOW\_rust\_verify\_mod.o := y}
  \item Boot kernel in virtme-ng: \texttt{vng --user root}
  \item Load modules: \texttt{insmod simple\_vuln\_mod.ko} and
        \texttt{insmod rust\_verify\_mod.ko}
  \item Run trigger tool for each case:
        \begin{itemize}[nosep]
          \item Opens \texttt{/dev/kcov\_dataflow}
          \item Calls \texttt{ioctl(KCOV\_DF\_INIT\_TRACE, 65536)}
          \item Maps buffer via \texttt{mmap()}
          \item Calls \texttt{ioctl(KCOV\_DF\_ENABLE)}
          \item Writes to procfs/debugfs trigger
          \item Calls \texttt{ioctl(KCOV\_DF\_DISABLE)}
          \item Parses and displays records from mapped buffer
        \end{itemize}
  \item Verify: compare captured values against \texttt{printk} ground truth
  \item Cross-check: verify KASAN report matches the corruption visible in
        ENTRY$\to$RET differential
  \item Symbol check: verify PC values in records match \texttt{/proc/kallsyms}
\end{enumerate}

\subsection{Kernel Symbol Verification}

After loading modules, verify kernel symbols are at expected addresses:
\begin{lstlisting}[numbers=none]
$ cat /proc/kallsyms | grep sanitizer_cov_trace
ffffffff88ab7600 T __sanitizer_cov_trace_args
ffffffff88ab78b0 T __sanitizer_cov_trace_ret
\end{lstlisting}

The PC values in kcov\_dataflow records should match module symbol addresses.
For modules, the base address is found in \texttt{/proc/modules}:
\begin{lstlisting}[numbers=none]
$ cat /proc/modules | grep simple_vuln
simple_vuln_mod 16384 0 - Live 0xffffffffc0400000
\end{lstlisting}

Records with \texttt{pc=0xffffffffc0400280} correspond to
\texttt{vuln\_process} at offset \texttt{0x280} from module base.

\subsection{drgn Verification of C Module}

Using drgn on \texttt{/proc/kcore} (live kernel) or a vmcore dump:
\begin{lstlisting}[language=Python,numbers=none]
import drgn
prog = drgn.Program()
prog.set_kernel()

# Verify symbol addresses match kcov_dataflow PCs
sym = prog.symbol("__sanitizer_cov_trace_args")
print(f"trace_args: 0x{sym.address:x}")
# Output: trace_args: 0xffffffff88ab7600

# For vmcore: verify C function locals
task = prog["init_task"]
print(f"kcov_df_size: {task.kcov_df_size.value_()}")
\end{lstlisting}

\subsection{drgn Verification of Rust Module (Failure Case)}

This demonstrates why \toolname{} is necessary for Rust:
\begin{lstlisting}[language=Python,numbers=none]
# After capturing vmcore during rust_process_data
prog = drgn.Program()
prog.set_core_dump("/tmp/vmcore-rust.elf")
prog.load_debug_info(["rust_verify_mod.ko"])

trace = prog.stack_trace(task)
for frame in trace:
    if "rust_process_data" in str(frame):
        print(f"Frame: {frame}")
        print(f"Locals: {frame.locals()}")
        # Output: Locals: {}
        # NO VARIABLES VISIBLE
        break
\end{lstlisting}

The empty locals dictionary confirms that rustc \texttt{-O2} elides all
\texttt{DW\_AT\_location} attributes for function parameters. The values exist
only in registers during execution and are not recoverable from a memory dump.

\subsection{Automated Regression Test Script}

We provide a complete test script that validates all five cases:
\begin{lstlisting}[language=bash,numbers=none]
#!/bin/bash
# test_kcov_dataflow.sh - automated verification

TRIGGER=/path/to/trigger
PASS=0; FAIL=0

run_test() {
  local name=$1 trigger=$2 expected=$3
  output=$($TRIGGER $trigger 2>&1)
  if echo "$output" | grep -q "$expected"; then
    echo "PASS: $name"
    ((PASS++))
  else
    echo "FAIL: $name (expected: $expected)"
    ((FAIL++))
  fi
}

# Load modules
insmod simple_vuln_mod.ko
insmod rust_verify_mod.ko

# Test cases
run_test "OOB" /proc/vuln_trigger "0x1337"
run_test "UAF" /proc/uaf_trigger "0x41414141"
run_test "DF"  /proc/df_trigger "0xdf00df00"
run_test "Rust" /sys/kernel/debug/rust_trigger \
  "0xcafe"

echo "Results: $PASS passed, $FAIL failed"
\end{lstlisting}

\section{Comparison with Alternative Approaches}\label{app:comparison}

\begin{table*}[t]
\centering
\caption{Comprehensive comparison of kernel data-flow observation approaches.}\label{tab:comparison}
\small
\begin{tabular}{@{}lcccccccc@{}}
\toprule
\textbf{Approach} & \textbf{Overhead} & \textbf{Annotation} & \textbf{Struct} & \textbf{Rust} & \textbf{No-crash} & \textbf{Temporal} & \textbf{Kernel} & \textbf{Always-on} \\
\midrule
KCOV (edge) & $<$1\% & None & \xmark & \cmark & \cmark & \xmark & \cmark & \cmark \\
ftrace/kprobes & 5--20\% & Per-func & \xmark & \xmark & \cmark & \cmark & \cmark & \xmark \\
eBPF & 2--10\% & Program & partial & \xmark & \cmark & \cmark & \cmark & \xmark \\
DFSan & 10--100$\times$ & None & \cmark & \xmark & \cmark & \cmark & \xmark & \xmark \\
KASAN & $<$5\% & None & \xmark & \cmark & \xmark & \xmark & \cmark & \cmark \\
drgn/vmcore & 0\% & None & \cmark & \xmark & \xmark & \xmark & \cmark & \xmark \\
\textbf{\toolname{}} & \textbf{$<$3\%} & \textbf{None} & \textbf{\cmark} & \textbf{\cmark} & \textbf{\cmark} & \textbf{\cmark} & \textbf{\cmark} & \textbf{\cmark} \\
\bottomrule
\end{tabular}
\end{table*}

Table~\ref{tab:comparison} provides a comprehensive comparison across eight
dimensions. \toolname{} is the only approach that satisfies all requirements
simultaneously: low overhead, zero annotation, struct field decomposition, Rust
support, no-crash detection, temporal ordering, kernel compatibility, and
always-on operation.

\paragraph{Key differentiators:}
\begin{itemize}[nosep]
  \item \textbf{vs.\ ftrace/kprobes:} Our approach is compile-time (zero
        runtime registration overhead) and provides struct field decomposition
        via debug metadata. kprobes require per-function setup and cannot
        automatically decompose struct arguments.
  \item \textbf{vs.\ eBPF:} eBPF programs are verified and restricted---they
        cannot perform arbitrary memory reads or access debug metadata at
        runtime. Our compile-time approach pre-computes field offsets.
  \item \textbf{vs.\ DFSan:} DFSan tracks byte-level taint propagation
        (complete data-flow graph) at enormous overhead. We capture only
        boundary snapshots (function entry/exit) at negligible cost---sufficient
        for fuzzing guidance and triage.
  \item \textbf{vs.\ drgn:} drgn requires a crash to produce a vmcore and
        fails on optimized Rust. We capture data continuously without crashes.
\end{itemize}

\section{Kernel Configuration Reference}\label{app:kconfig}

\begin{lstlisting}[numbers=none]
# Required base
CONFIG_KCOV=y
CONFIG_DEBUG_INFO_DWARF5=y

# Dataflow tracing (our addition)
CONFIG_KCOV_DATAFLOW_ARGS=y
CONFIG_KCOV_DATAFLOW_RET=y

# Recommended for full verification
CONFIG_KASAN=y
CONFIG_KASAN_INLINE=y
CONFIG_RUST=y

# Module Makefile (single line to enable)
KCOV_DATAFLOW_my_module.o := y
\end{lstlisting}

\subsection{Kconfig Dependencies}

\begin{lstlisting}[numbers=none]
config KCOV_DATAFLOW_ARGS
    bool "Enable dataflow tracing of function arguments"
    depends on KCOV
    help
      Enables compile-time instrumentation that captures
      function arguments at entry, including automatic
      struct field expansion via DWARF debug metadata.

config KCOV_DATAFLOW_RET
    bool "Enable dataflow tracing of return values"
    depends on KCOV
    help
      Enables compile-time instrumentation that captures
      function return values, including struct field
      expansion for composite return types.
\end{lstlisting}

\section{Raw Evaluation Output}\label{app:raw}

\subsection{C Module: OOB Write (Full Trace)}

\begin{lstlisting}[numbers=none]
=== KCOV Dataflow TLV Dump (50 words) ===
[ENTRY] seq=1 pc=0xffffffffc0400010 arg[0](8)
  ptr=0x888006578300
  val = 0x0
  val = 0x40d801a
  val = 0xffffffffc02a0000
[ENTRY] seq=2 pc=0xffffffffc0400010 arg[1](8)
  ptr=0x480132
  val = 0x0
[ENTRY] seq=3 pc=0xffffffffc0400010 arg[2](8)
  ptr=0x8880047e7d70
  val = 0x1
[ENTRY] seq=4 pc=0xffffffffc0400010 arg[3](8)
  ptr=0x8880047e7ec8
  val = 0x0
[ENTRY] seq=5 pc=0xffffffffc0400280 arg[0](8)
  ptr=0x888001f09d00
  struct session_data (3 fields):
    val = 0x1337        # .id
    val = 0x696e697469  # .buf ("initi")
    val = 0xf           # .size = 15
[RET]   seq=6 pc=0xffffffffc0400280 ret(8)
  ptr=0x888001f09d00
  struct session_data (3 fields):
    val = 0x1337        # .id (unchanged)
    val = 0x0           # .buf (overwritten)
    val = 0x20          # .size = 32 [CORRUPT]
[RET]   seq=7 pc=0xffffffffc0400010 ret(8)
  ptr=0x8880047e7d68
  val = 0x1
=== Done ===
\end{lstlisting}

\subsection{Rust Module: Full Trace (Abbreviated)}

\begin{lstlisting}[numbers=none]
=== KCOV Dataflow TLV Dump (372 words) ===
[ENTRY] seq=9 pc=0xffffffffc04028c0 arg[0](8)
  ptr=0x888007f6fdb0
  struct RustData (3 fields):
    val = 0xcafe              # .id
    val = 0xdeadbeef12345678  # .value
    val = 0xa                 # .flags
[ENTRY] seq=10 pc=0xffffffffc04028c0 arg[1](4)
  ptr=0x888007f6fd44
    val = 0x3                 # multiplier
[ENTRY] seq=11 pc=0xffffffffc04028c0 arg[2](8)
  ptr=0x888007f6fd08
    val = 0x1111111111111111  # cookie
... (internal function calls: seq 12-74) ...
[RET] seq=75 pc=0xffffffffc04028c0 ret(8)
  ptr=0x888007f6fca0
    val = 0xefbed00023456789  # return value
=== Done ===
\end{lstlisting}

\subsection{KASAN Report (OOB Case)}

\begin{lstlisting}[numbers=none]
[    4.509361] BUG: KASAN: slab-out-of-bounds in
  vuln_process+0x43/0xb0 [simple_vuln_mod]
[    4.509365] Write of size 32 at addr
  ffff888001f09d10 by task trigger/234
[    4.509370] CPU: 2 PID: 234 Comm: trigger
[    4.509375] Call Trace:
[    4.509378]  dump_stack_lvl+0x5d/0x80
[    4.509382]  kasan_report+0xd8/0x110
[    4.509386]  vuln_process+0x43/0xb0
[    4.509390]  vuln_write+0x89/0xc0
[    4.509394]  proc_reg_write+0x12e/0x1a0
\end{lstlisting}

Note: KASAN reports the address (\texttt{ffff888001f09d10}) and size (32) but
not which struct field was corrupted or what the original value was. \toolname{}
provides both: the field is \texttt{.size} (offset 12), original value was 15,
corrupted to 32.

\end{document}
